\chapter{\rm{VASP}软件的理论基础}
\textrm{VASP}是维也纳大学的\textrm{George Kresse}团队开发和维护的一款基于第一原理的电子结构计算和分子动力学模拟软件包。
在常用的微观尺度计算材料软件中，\textrm{VASP}是公认的执牛耳者。
总体来说，\textrm{VASP}在同类软件中能脱颖而出，主要得益于其在以下方面的性能优势:
\begin{itemize}
	\item \textrm{VASP}采用\textrm{PAW~(Projector Augmented-Wave)}方法，平衡了赝势方法和全电子计算优点，兼顾了计算的精度和效率。
		\item \textrm{VASP}在实空间优化投影函数\textrm{(Projector)}，将主要的计算过程变换到实空间完成，节省了矩阵迭代的内存开销。%，保证了计算精度和效率
		\item \textrm{VASP}通过引入多样的优化算法，提高了矩阵对角化和电荷密度搜索的效率。
		\item 在\textrm{VASP}的并行计算中，有效均衡了各节点处理\textrm{FFT}变换负载和通信，提升了软件的并行效率。近年来，在\textrm{NVIDA}的支持下,\textrm{VASP}的\textrm{GPU}版，计算效率又有了很大的提升。
\end{itemize}
和很多科学计算软件的发展过程类似，\textrm{VASP}软件与先期发展的同类型软件也经历了若干版权纷争的曲折，最后以独立的商业软件形式面世。
%\textrm{VASP}软件诞生之前，
大约在1980年代中后期，当时在麻省理工学院\textrm{(Massachusetts Institute of Technology，MIT)}做访问研究工作的\textrm{Mike Payne}开发了一套基于密度泛函理论和赝势的从头计算程序，\textrm{Mike Payne}回到英国剑桥大学后，将该程序完善成为\textrm{CASTEP(Cambridge Sequential Total Energy Package)}软件包。
1989年7月，\textrm{J\"urgen Hafner}将在剑桥得到的\textrm{Mike Payne}组的代码带回维也纳%，由此开启\textrm{VASP}软件的``初始''代码开发时期；
并于1991年9月，在\textrm{CASTEP-1989}年版本基础上，正式开启新的软件代码编写工作。与此同时，\textrm{CASTEP}代码也在作进一步开发。1992年10月起，代码中开始引入超软赝势\textrm{(Ultra-Soft Pseudo-Potential, USPP)}；1995年2月，软件正式定名\textrm{VASP}(一度曾命名为\textrm{VAMP})；1996年9月，自\textrm{VASP}-4.1版起，用\textrm{Fortran~90}语言改写全部代码，并于1997年1月前后完成\textrm{MPI}并行版本的开发，有必要指出的是。此时\textrm{VASP}团队的并行版代码中，并行通信实现的核心实现同样源自\textrm{CASTEP}，这一阶段的竞争，引发了两个团队间的龃龉。
1998年7月前后，\textrm{VASP}团队不得不彻底改写并行的通信实现，完全替代\textrm{CASTEP}的实现方案\footnote{正因此，\textrm{VASP}组坚称，\textrm{J. Hafner}带回的代码虽与\textrm{CASTEP}同源，但彼时\textrm{CASTEP}名份未定，因此属于两者各属不同的分枝；后来\textrm{VASP}坚决发展\textrm{PAW}方法，与\textrm{CASTEP}划清界限，皆源于此。科学软件发展史上类似事件层出不穷，在此不作专门介绍和评论。}
1999年12月，\textrm{VASP}推出应用引入投影子缀假波\textrm{(Projector Augmented  Wave, PAW)}方法的全新版\textrm{VASP}，从此与\textrm{CASTEP}彻底分道扬镳。
2004年起，\textrm{VASP}陆续推出基于\textrm{PAW}的稳定版本\textrm{5.X}，2015年前后，\textrm{VASP}实现了对GPU的支持，并经历了从\textrm{CUDA-C}到\textrm{OpenACC}的技术演进，从2015年底的\textrm{VASP-5.4.1}开始，通过\textrm{CUDA-C}编程，支持\textrm{VASP}的\textrm{GPU}计算；2021年推出\textrm{VASP-6.2.0}，主推\textrm{OpenACC}端口，极大提升了软件的速度与代码可维护性，正式开启支持\textrm{GPU}计算的\textrm{VASP}稳定版本发布。
截止到目前，最新版本的\textrm{VASP-6.5.0}已经实现了\textrm{(Bethe-Salpeter Equation, BSE)}等更多功能的\textrm{GPU}加速。需要说明的是，当前的\textrm{VASP}仅支持\textrm{NVIDIA}的\textrm{GPU}，\textrm{AMD}与\textrm{Intel}的\textrm{GPU}支
持计划仍在开发中。

以下对\textrm{VASP}软件所涉及的密度泛函和第一原理分子动力学相关的基础理论作扼要介绍，以便于读者整体了解\textrm{VASP}的软件框架和底层逻辑。

\section{\rm{VASP}软件的基础}
\textrm{VASP}软件的理论基础主要包括密度泛函理论，电子能带结构理论，具体的计算方法则涉及投影子缀加波方法，$\vec k$空间积分和总能量的计算，\textrm{DFT}+$U$的处理等技术方案的实现。
\subsection{密度泛函{理论}} \label{chap:dft}
密度泛函理论\upcite{Parr-Yang,CR91-651_1991}%最早})
可溯源到二十世纪二十年代的Thomas-Fermi%理论})
{模型}\upcite{PCPS23-542_1927,ZP48-73_1928}和五十年代Slater提出的X\,$\alpha$方法\upcite{PR81-385_1951}。经典量子%化})
{力}学的基本变量是波函数，而密度泛函理论的%最基本的变量是体系单粒子})
{基本变量是体系的电子}密度，通过体系的%单})
粒%{电}
子密度确定体系基态和性质，如总能量等，降低了计算的复杂度。直到六十年代，Hohenberg和Kohn\upcite{PR136-B864_1964}严格%地})
证明了多%粒})
{电}子%量子力学})
体系的%基态
性质%可以})
由基态%单})
%粒})
{电}子密度完全决定，且体系精确的基态%单})
%粒})
{电}子密度%使得})
{对应于}基态总能量%变分取极小})
{泛函的最小极小值}，为密度泛函理论奠定了坚实的理论基础。%因为})
然而，体系的基态总能量%和})
和基态%单})
%粒})
{电}子密度%之间的具体对应关系并不清楚，特别是体系的动能与基态单粒子的对应关系对于精确计算体系基态总能量有比较重要的影响})
{的泛函的显式表达}至今未知。Kohn和\linebreak Sham\upcite{PR140-A1133_1965}%采用了与相互作用体系有相同单粒子密度的无相互作用体系的动能近似来近似计算相互作用体系的动在在能，把剩余部分都放在交换-相关能项中予以考虑。本质上，这个方法把单粒子波函数又引入了密度泛函理论，但是基本解决了动能计算的问题。在他们的工作基础上，各种近似的交换-相关能泛函形式\upcite{CJP58-1200_1980,PRB45-13244_1992,PRA38-3098_1988,PRB33-8822_1986,PRB37-785_1988}被提出来并})
{提出能量近似计算方案，使密度泛函理论可应用于实际体系的计算。后来人们提出过多种能量密度泛函的近似表达式，其中有一部分已被}广泛应用于%计算各类})
原子、分子和固体体系{的计算中\upcite{CJP58-1200_1980,PRB45-13244_1992,PRA38-3098_1988,PRB33-8822_1986,PRB37-785_1988}}。

\subsubsection{Hohenberg-Kohn定理}
Hohenberg-Kohn定理\upcite{PR136-B864_1964}包含两部分{内容}。第一部分指出，多%粒})
{电}子体系的外势由其%基态性质与
基态%单})
%粒})
{电}子密度%之间存在一一对应关系，这部分是密度泛函理论的基础
决定(除可能相差一个常数以外)，从而体系的基态波函数以及所有其它性质均由其基态电子密度决定
；第二部分指出体系基态总能量{泛函}在体系基态%单})
%粒})
{电}子密度处取极小{值}，这部分%给出})
{提供}了从%单}) %粒})
{电}子密度计算体系基态能量的基本方法。Hohenberg-Kohn定理的证明如下：

%设})
N个粒子相互作用量子力学体系的%Hamiltonian})
{基态总能量}为：
\begin{equation}
	{E_{\mathrm{gs}}[\Psi]}%F_{\mathrm{HK}}[\rho]})
=\langle\Psi%[\rho]})
|\hat{T}+\hat{V}+\hat{W}|\Psi%[\rho]})
\rangle
	\label{eq:DFT_01}
\end{equation}
其中$\hat{T}$是动能算符，$\hat{V}$是外势%、即通常的局域单粒子势})
，$\hat{W}$为%粒})
{电}子间相互作用势{，$\Psi$为基态波函数}。假设两个体系的外势分别为$\hat{V}$和$\hat V^{\prime}$，它们之间相差不为常数，这两个体系的非简并基态分别为$\Psi$和$\Psi^{\prime}$，其基态的本征方程分别为：
$${\hat H\Psi=}\bigl(\hat{T}+\hat{V}+\hat{W}\bigr)|\Psi\rangle=E_{\mathrm{gs}}|\Psi\rangle$$
$${\hat H^{\prime}\Psi=}\bigl(\hat{T}+\hat V^{\prime}+\hat{W}\bigr)|\Psi^{\prime}\rangle=E_{\mathrm{gs}^{\prime}}|\Psi^{\prime}\rangle$$

%假设})
{根据量子力学原理，}$|\Psi\rangle{\neq}|\Psi^{\prime}\rangle$，%则上两式相减可得})：
% $$%\big(\hat{V}-\hat V^{\prime}\big)|\Psi\rangle=\big(E_{gs}-E^{\prime}_{gs}\big)|\Psi\rangle}) $$
%因为$\hat{V}$和$\hat V^{\prime}$是乘积算符，从上式可知，$\big(\hat{V}-\hat V^{\prime}\big)=\big(E_{gs}-E^{\prime}_{gs}\big)=\text{const}$，这与$\hat{V}$和$\hat V^{\prime}$之间相差不为常数矛盾，因此$\hat{V}$和$|\Psi\rangle$之间存在一一对应关系。})
由变分原理可知：
$$E_{\mathrm{gs}}=\langle\Psi|\hat{H}|\Psi\rangle<\langle\Psi^{\prime}|\hat{H}|\Psi^{\prime}\rangle$$

设$|\Psi\rangle$和$|\Psi^{\prime}\rangle$对应的%单})
%粒})
{电}子密度分别为$\rho(\vec{r})
$和$\rho^{\prime}(\vec{r})$，{$v(\vec r)$为单个电子受到的外势的作用能，则}有
$$\langle\Psi^{\prime}|\hat{H}|\Psi^{\prime}\rangle=\langle\Psi^{\prime}|\hat H^{\prime}+\hat{V}-\hat V^{\prime}|\Psi^{\prime}\rangle=E^{\prime}_{gs}+\int{\rho^{\prime}(\vec{r})[v(\vec{r})-v^{\prime}(\vec{r})]\textrm{d}^3\vec{r}}$$
%相应的，从$E^{^{\prime}}_{gs}$出发})
{类似地}可以得到不等式
$$E^{^{\prime}}_{gs}<E_{gs}+\int{\rho(\vec{r})[v^{\prime}(\vec{r})-v(\vec{r})]\textrm{d}^{3}\vec{r}}$$

%假设})
{如果}$\rho(\vec{r})\!=\!\rho^{\prime}(\vec{r})$，把上面两个不等式相加可以得到$E_{gs}+E^{\prime}_{gs}\!<\!E_{gs}+E^{\prime}_{gs}$，这显然%与假设矛盾，因此$\rho(\vec{r})$与v可以表示的基态$|\Psi\rangle$之间有一一对应关系})
{不能成立。亦即不可能出现不同的外势产生相同的电子密度的情况}。由此可以得出结论，体系的基态%单})
%粒})
{电}子密度和体系的外势(除了相差一个常数之外)一一对应%，})
{。}%也就是说，})
{因$\rho(\vec r)$唯一地确定体系中电子的数目$N$，故}体系%的一切})
{%基态
的\textrm{Hamiltonian}、从而它的全部}性质都%可以})
由体系的基态%单})
%粒})
{电}子密度完全决定。{特殊地，体系基态总能量是它的基态电子密度的泛函，记作$E[\rho]$。
	
设有一电子密度$\tilde\rho(\vec r)$，满足约束条件
\begin{displaymath}
	\begin{aligned}
		&\tilde\rho(\vec r)\!\geqslant\!0\\
		&\int\tilde\rho(\vec r)\textrm{d}^3\vec r\!=\!N
	\end{aligned}
\end{displaymath}
$\tilde\rho(\vec r)$来自某个\textit{v}\,-\,可表示的基态波函数$\tilde\Psi$(称为\textit{v}\,-\,可表示的电子密度)。设与体系基态波函数$\Psi$对应的电子密度为$\rho_0(\vec r)$，若$\tilde\rho(\vec r)\!\neq\!\rho_0(\vec r)$，则根据上面的讨论可知，$\tilde\Psi\!\neq\!\Psi$，$E[\tilde\rho]$\,$\geqslant$\,$E[\rho_0]$。}

%从上面的结论我们可以定义作为单粒子密度$\rho(\vec{r}) $泛函的体系基态能量的形式。设一$N$个粒子体系的外势场为$\hat{V_{0}}$，})
{设}与$\rho(\vec{r})$%一一})
对应的%某一N个粒子})
体系波函数为$|\Psi[\rho]\rangle$，%定义该})
体系的基态总能量作为$\rho(\vec{r})
$的泛函%形式})
{可表达}为：
\begin{equation}
	E[\rho]=\langle\Psi[\rho]|\hat{T}+\hat{V}+\hat{W}|\Psi[\rho]\rangle
	\label{eq:DFT_02}
\end{equation}
%由变分原理可知，体系本身的基态单粒子密度使得$E_{v_{0}}[\rho]$取极小，因此可以通过寻找能量泛函的极小值点来确定体系的基态能量和基态粒子密度}
{因为由外势产生的势能是电子密度$\rho(\vec r)$的简单泛函，}体系的{基态}总能量泛函可以写成：
\begin{equation}
	E[\rho]=F_{\mathrm{HK}}[\rho]+\int\rho(\vec{r})v(\vec{r})\textrm{d}^3\vec{r}
	\label{eq:DFT_03}
\end{equation}
其中$F_{\mathrm{HK}}[\rho]=\underset{\Psi\to\rho}{\mathrm{Min}}\langle\Psi[\rho]|\hat{T}+\hat{W}|\Psi[\rho]\rangle$%。})
{，}%从上式可以看出，$F_{\mathrm{HK}}$的形式}
与%$v_{0}$})
{$v(\vec r)$}没有关系，%只与动能和离子间相互作用能的形式有关。对于我们最通常研究的原子、分子和固体体系而言，$F_{\mathrm{HK}}$的形式都是一致的。但是})
{是一个普适的能量泛函。}

到目前为止，$F_{\mathrm{HK}}$的%具体})
{精确表达形}
式还没有得到，%一般都采用近似的方法})
{只能作近似处理}。Thosmas-Fermi%理论})
{模型}%实际上就是一个对$F_{\mathrm{HK}}$近似的理论，但是由于动能泛函的误差很})
{可以看作是对$F_{\mathrm{HK}}[\rho]$泛函采取粗略近似的密度泛函理论。由于误差太}大，Thomas-Fermi%理论})
{模型}既不能%得到})
{说明}原子体系%正确的})
电子密度的壳层结构，也不能%得到})
{说明存在}稳定%的})
分子{的事实}。%后来经过})
Dirac\upcite{PCPS26-376_1930}和Weizsacker\upcite{ZP96-431_1935}对能量泛函的改进取得%了})
一定%的进展})
{成功}，但%是同实验结果相比，仍然有比较大的差距})
{没有取得实质性的进展}。

上述%的证明})
只讨论了非简并基态的问题{。}Kohn证明\upcite{Bassani-Fumi-Tosi}，对于%简并体系基态})
{体系简并基态}的情况，虽然%单})
%粒})
{电}子密度和体系的基态波函数之间的一一对应关系不再成立，但是同样可以定义一个普适的能量泛函$F_{\mathrm{HK}}$，体系的基态%单})
%粒})
{电}子密度使%得})
体系%的})
基态{总}能量{泛函}取极小{值的定理}仍然成立。%前})
{在上}面的证明%隐含的一个})
{中}假设%是我们所取的})
%单})
%粒})
{电}子密度$\rho$是\textit{v}\,-\,可表示的，%即$\rho$是由某个外势决定的Hamiltonian算符的基态单粒电子密度})
{这使得通过变分方法寻找能量泛函的极小值难于实现}。Levy\upcite{PNAS76-6062_1979}证明，实际上，只要$\rho$是
%\it v}-})
{$N$}可表示的，即$\rho$%是与})
{来自}某个$N$电子体系的基态波函数%相对应})
，上面的定理也%能够})
成立。可以证明\upcite{PRB12-2111_1975,IJQC24-243_1983}，任意一个合理的%单})
%粒})
{电}子密度都是%{\it v}-}
{$N$}可表示的，这样就%可以})
基本解决{了}变分空间的%选取})
问题。

\subsection{Kohn-Sham方程}
%对于})
{根据}Hohenberg-Kohn定理，只要%得到})
{知道}体系的总能量的{密度}泛函形式，就可以通过对%单})
%粒})
{电}子密度%作})
变分求得体系的基态%单粒})
{总能量和基态的电}子密度%和基态能量})
。与%通常的})
对波函数%作})
变分求%得})
体系基态%性质})
{总能量}的方法相比，波函数是3N%-维})
{维空间中的}函数，而%单})
%粒})
{电}子密度只是三维{空间中的}函数，因此这种方法有%巨})
{很}大的优越性。但是%由于很难})
{迄今无法}得到%比较})
精确的能量密度
泛函形式，特别是体系的动能密度
泛函的精确%度对计算体系性质有很大的影响})
{形式。为避开这一困难}，Kohn和Sham\upcite{PR140-A1133_1965}提出%了用具有与相互作用体系单粒电子密度相同但无相互作用体系的动能来近似计算相互作用体系动能的方法，这种})
{以下}方法%比较好地})
解决%了})
动能泛函的计算问题。%他们的})
具体%计算})
方法如下：假设存在一个外势为$v_s$的无相互作用多%粒})
{电}子体系，其基态%单})
%粒})
{电}子密度$\rho$与%我们所})
{要}研究的一个多电子相互作用体系的基态%单})
%粒})
{电}子密度相同。这个无相互作用体系的%单})
%粒})
{电}子密度为：
$$\rho(\vec{r})=\sum_{i=1}^{N}\langle\varphi_{i}|\vec{r}\rangle\langle\vec{r}|\varphi_{i}\rangle$$
其中$\varphi_i(i\!=\!1,\cdots,N)$是该体系能量最低的$N$个{单电子}本征态。该体系的动能为：\linebreak $T_{s}[\rho]\!=\!\sum\limits_{i=1}^{N}\langle\varphi_{i}|\hat{T}|\varphi_{i}\rangle$。
用这个无相互作用体系的动能作为%我们})
{要}研究的相互作用体系的动能近似{值}，可以把相互作用体系的%动能})
{基态能量密度泛函}写成如下形式：
\begin{equation}
  \label{eq:dft-1}
E[\rho]=T_{s}[\rho]+\int{v(\vec{r})\rho(\vec{r})\textrm{d}^3\vec{r}}+\frac 12\iint{\rho(\vec{r})w(\vec{r},\vec r\,^{\prime})\rho(\vec r\,^{\prime})\textrm{d}^3\vec{r}\textrm{d}^3\vec r\,^{\prime}}+E_{XC}[\rho]
\end{equation}
其中{$w(\vec r, \vec r\,^{\prime})$为两个电子间的相互作用能。}交换-相关能泛函的形式为：

\begin{equation} 
  \label{eq:dft-2}
  E_{XC}[\rho]=F_{\mathrm{HK}}[\rho]-\frac 12\iint{\rho(\vec{r}) w(\vec{r},\vec r\,^{\prime})\rho(\vec r\,^{\prime})\textrm{d}^3\vec{r}\textrm{d}^3\vec r\,^{\prime}}-T_{s}[\rho]
\end{equation}
利用$T_s$和$\rho$的表达式，%对})
{由}总能量%作})
{在$\{\varphi_i, i\!=\!1,\cdots, N\}$满足正交归一条件下对$\{\varphi_i\}$}变分，即%可以})
得到Kohn-Sham方程：
\begin{equation} \label{eq:dft-3}
	\bigl(\hat{T}+\hat{V}_{\mathrm{eff}}\bigr)|\varphi_{i}\rangle=\varepsilon_{i}|\varphi_{i}\rangle{,\quad i=1,\cdots,N,\cdots}
\end{equation}
其中
\begin{equation} \label{eq:dft-4}
	V_{\mathrm{eff}}(\vec{r})={v}(\vec{r})+\int{w(\vec{r},\vec r\,^{\prime})\rho(\vec r\,^{\prime})\textrm{d}^3{\vec r\,^{\prime}}+\frac{\delta E_{XC}}{\delta\rho(\vec{r})
}}
\end{equation}
式\eqref{eq:dft-4}右边第一项%中})
${v}(\vec{r})
$为外势，第二项为%粒})
{电}子间经典的{Coulomb}相互作用势，第三项是交换-相关势。%对于我们所关心的原子、分子和固体等多电子体系而言，\eqref{eq:dft-4}式中的外势为核吸引势，粒子间经典相互作用势是电子间的Coulomb势。})

以上推导中假设任意一个相互作用%的{\it v}-可表示的单粒})
{体系的基态电}子密度%也是})
{有一个}无相互作用%的{\it v}-可表示单粒})
{体系的基态电}子密度{与之相同}。可以证明\upcite{PNAS76-6062_1979,IJQC24-243_1983}，按照该方法构造的无相互作用体系基态如果是非简并的，即这个无相互作用体系的基态可以用单行列式波函数表示，则用这个单行列式波函数构造的%单粒})
{基态电}子密度与相互作用体系的%单粒})
{基态电}子密度是一致的%，但是如果无相互作用体系的基态是简并的，相互作用{\it v}-可表示单粒子密度是否仍然与由单行列式波函数构造的无相互作用体系单粒子密度对应的问题还没有得到解决})
。

对于{电子}自旋{密度}%极化的情况，从})
{不恒为零的体系，推广}Hohenberg-Kohn定理%可以知道，体系基态自旋为})
{的推导过程和近似处理方法，可以得到对于体系的}$\alpha$%})
自旋和%为$\beta$的})
{$\beta$自旋}电子的{\textrm{Kohn-Sham}方程，被称为自旋密度泛函理论\textrm{(SDFT)}。}%密度从原理上可以由体系总电子密度完全决定；但是如果要研究})
{在}研究一些体系的磁学性质或者处于解离极限的分子的性质{时}，%在现有泛函基础上，只有})
采用%把自旋$\alpha$和自旋$\beta$的电子分开处理的})
自旋密度泛函方法才能得到合理的结果{。}%而且})
对于开壳层体系%而言})
，用自旋密度泛函方法%的结果})
通常比不{区}分自旋的密度泛函方法结果要好。%自旋密度泛函理论的基本原理和Kohn-Sham方程很容易用类似的推导方法得到，在此不作推导了。})

Kohn-Sham方程的%这个})
实质，是把动能泛函的主要部分分离出来，将剩余%不太重要的动能})
部分放到交换-相关能中。由于该方程%基本})
{比较好地}解决了%计算})
动能泛函的{计算}问题，%这个方法})
使得DFT能够比较精确地研究电子体系的性质。有必要指出，Kohn-Sham方程虽然把单%粒})
{电}子波函数引进密度泛函理论，但它是{形式上}的单粒子方程，%从理论上讲只利用单Slater行列式就})
可以通过简单的交换-相关泛函%计算})
{，近似地考虑}电子间的多体关联，%而传统的量子化学方法如果要考虑电子之间的相关能效应，则必须采用多Slater行列式的方法，Kohn-Sham方程})
{相比}比传统的量子%化})
{力}学方法，有明显的%})
优越%性})
{之处}，{因此}基于Kohn-Sham方程的密度泛函方法也成为电子尺度计算最重要的{理论}方法之一。

\subsubsection{交换-相关能泛函}
Kohn-Sham方程%从})
理论上%讲})
是%一个})
严格的方程，但是要把它用于具体计算还需要%解决})
{知道}交换-相关能泛函的具体形式。%从上一节的推导可以看出，})
Kohn-Sham%理论})
提出的交换-相关能的含义与从头计算方法的交换和相关能的含义有%很})
{较}大的差别{。现有}密度泛函理论中的交换能定义为用KS轨道按照Hartree-Fock交换能%形})
{表达}式计算得到的能量，相关能则包含两部分，一部分是相互作用体系动能和对应的无相互作用体系动能{之}差，另一部分是相互作用体系中粒子间的相互作用能与经典相互作用{能}和交换能%的})
{之}差。密度泛函{理论}中的交换-相关能泛函的精确形式还%没有})
{无法}得到，只能通过{考虑}一些{应满足的}具体%的})
物理%考虑})
{条件}和合理的假设得到近似的交换-相关能泛函。最简单的近似就是由Kohn和Sham提出的局域密度近似%方法})
{泛函}\upcite{PR140-A1133_1965}(local density approxmiation, LDA){。}局域密度近似基于均匀电子气模型，一般认为对于电子密度变化缓慢的体系{才会}是%个})
比较好的近似，而实际上对于原子、分子和某些价电子密度%变化})
剧烈变化的固体，LDA方法仍然能够得到比较合理的结果。不过，为了进一步提高密度泛函方法对原子、分子以及扩展体系的计算精度，需要考虑%到})
电子密度的不均匀性。通过引入电子密度的梯度对{定域近似}交换-相关能{泛函}的形式加以校正，%即所谓的})
{就是}梯度展开近似(Gradient Expansion Approximation, GEA)%和})
{或}一般梯度近似(Generalized Gradient Approximation, GGA)，%即})
能量泛函不仅依赖于局域密度值的分布，还与密度梯度有关，这在一定程度上改善LDA的计算结果。比较重要的GGA交换能泛函有P86$_x$\upcite{PRB33-8800_1986}、B88\upcite{PRA38-3098_1988}、PW91\upcite{PRB46-6671_1992,PRB48-4978_1993,PRB54-16533_1996,PRB57-14999_1998}、PBE\upcite{PRL77-1396_1996,IBID78-1396_1997}、HCTH\upcite{JCP109-6264_1998}等。

由于在交换-相关能泛函的定义中包含了一部分动能的贡献，因此很多研究者考虑{到}在交换-相关{能}泛函中应该包括动能密度为变量，%这就是所谓的})
{提出}\textit{meta}-GGA泛函。自1980年代末以来，比较有代表性的\textit{meta}-GGA交换能泛函有BR89\upcite{PRA39-3761_1989}、\linebreak VSXC\upcite{JCP109-400_1998}、tauPBE\upcite{JCP111-911_1999}、PKZB\upcite{PRL82-2544_1999,PRL82-5179_1999}、B00\upcite{JCP112-4020_2000}、PCS00\upcite{JCP113-10013_2000}等。

为了提高各种近似{能量}泛函的%计算结果的})
精度，从1990年代初起，基于绝热关联%(下面会详述)})
{的推理}，Becke提出把部分按Hartree-Fock方法计算的交换能加入到能量密度泛函中的思想{，}%从而})
构造出了一系列%的所谓})
杂化型%的})
交换-相关{能}泛函，其中比较重要的是“半对半(half and half)”泛函\upcite{JCP98-1372_1993}、B3P\upcite{JCP98-5648_1993}、B3LYP\upcite{JPC98-11623_1994}、B1B95\upcite{JCP104-1040_1995}等。

由于%对})
相关能的物理%定})
{含}义不%很})
{够}清晰，而且相关能的数值比交换能小一个数量级，既比较难准确估计，又不是计算中误差的最主要来源，因此提出的{近似}相关能泛函%模型})
比交换能泛函%模型})
少一些。最早提出把相关能表示成电子密度泛函的是Wigner\upcite{PR46-1002_1934}，他的公式适用于均匀电子气：
\begin{equation}
  E_C=-\dfrac a{r_s+b}
  \label{eq:dft-16}
\end{equation}
其中$a$、$b$是参数，$r_s$\,=\,$\left(\dfrac3{4\pi\rho}\right)^{\frac13}$称为表观半径。%实际上})
对%于})
均匀电子气的相关能并没有统一\linebreak 的解析式。von Barth\upcite{JPC5-1629_1972}和Gunnarsson\upcite{PRB13-4274_1976}等人都提出过一些近似的式子。1980年，\linebreak Vosko、Wilk和Nusair\upcite{CJP58-1200_1980}利用Ceperley和Alder\upcite{PRL45-566_1980}用量子Monte Carlo方法计算得到的{均匀电子气}相关能数据拟合出一个%均匀电子气的})
相关能的%公})
{表达}式，即最常用的VWN%公})
{表达}式。后来Perdew和Wang\upcite{PRB45-13244_1992}把他们的%公})
{表达}式简化并重新拟合了参数。

以上相关能%公式})
{泛函}都是针对均匀电子气的，比较常用的考虑了梯度校正的相关能泛函%公})
{表达}式有P86$_c$\upcite{PRB33-8822_1986}、LYP\upcite{PRB37-785_1988}、Bc88\upcite{JCP88-1053_1988}、PW91\upcite{PRB46-6671_1992,PRB48-4978_1993,PRB54-16533_1996,PRB57-14999_1998}等。

按照Perdew的建议\upcite{Perdew-Schmidt_2001}，现有的交换-相关能密度泛函可以分为以下几类：
\begin{enumerate}
  \item LDA: 即泛函只与密度分布的局域值有关；
  \item GGA: 泛函所依赖的变量除了局域密度之外，还包括局域密度的梯度；
  \item \textit{meta}-GGA: 泛函依赖的变量还包括动能密度；
  \item {杂化(hybrid)泛函}: 泛函与占据轨道有关；
  \item 完全非局域泛函: 泛函与所有的占据和非占据轨道都有关，属于理想泛函但不%现实})
{切合实际}。
\end{enumerate}
\begin{figure}[!h]
\centering
\includegraphics[height=3.4in,width=3.18in,viewport=10 5 680 700,clip]{Jacobi-ladder.png}
\caption{\small The schematic digram of Jacob's Ladder for DFT.\upcite{Perdew-Schmidt_2001,Science298-759_2002}}
\label{Fig:Jacob-Ladder}
\end{figure}

上述{各类}泛函%的类别})
从上到下越来越接近%化学})
精确值，因此被形象地视为``Jacob's Ladder''(登天梯，见图\ref{Fig:Jacob-Ladder}):~将\textrm{LDA~}视为第一级阶梯，\textrm{GGA~}视为第二级阶梯，更复杂的泛函形式视为第三级及更高的阶梯。\upcite{PRL91-146401_2003}但在密度泛函中过度引入轨道将会%造成})
{大大增加}计算量%的大大增加})
，失去密度泛函理论对%一般})
从头计算方法的优势，因此是{并}不可取的。
下面对每一类中最重要的泛函作一些简单的介绍。

\subsubsection{局域密度近似和LDA泛函}
局域密度近似是Kohn和Sham提出的最简单的一种近似处理交换-相关能的方法。交换-相关能{泛函}的具体形式取为：
\begin{equation}
  \label{eq:dft-5}
  E_{XC}^{\mathrm{LDA}}[\rho]=E_X^{\mathrm{LDA}}[\rho]+E_C^{\mathrm{LDA}}[\rho]=\int\varepsilon_X[\rho]\rho(\vec{r}) \textrm{d}^3\vec{r}+\int\varepsilon_C[\rho]\rho(\vec{r}) \textrm{d}^3\vec{r}
\end{equation}
其中$\varepsilon_X[\rho]$和$\varepsilon_C[\rho]$分别为{一个电子的}交换能和相关能%密度})
。上式表明空间每一个点的交换能密度和相关能密度只是取决于该点的电子密度，而与其他点的电子密度无关。对{于均匀电子气，}$\varepsilon_X[\rho]$和$\varepsilon_C[\rho]$%的选取可按照均匀电子气中交换-相关能和电子密度的关系得到，})
{为常数%，等于体系的总电子交换能(或相关能)除以体系的体积和电子密度
。}对于非均匀电子气%而言})
，%这样选取的$e_X(\rho)$和$e_C(\rho)$的物理图象相当于})
{可以设想}把空间分割为无穷多个%非常})
{无限}小的区域，%假设})
{可以认为}电子的密度在每个小的区域内都是均匀的，每个小区域中的交换-相关能都%采用})
{可以按}均匀电子气%模型下})
的交换-相关能{计算}，整个体系的交换-相关能就可以写成\eqref{eq:dft-5}式的形式。{采用}这种模型{，}%下的})
\textrm{Kohn-Sham~}方程的交换势表达为：
\begin{equation} \label{eq:dft-6}
	V_X^{\mathrm{LDA}}(\vec{r}) =\dfrac {\delta E_X^{\mathrm{LDA}}}{\delta\rho(\vec{r}) }=-\dfrac 32\alpha\left\{\dfrac3\pi\rho(\vec{r}) \right\}^{1/3}
\end{equation}
Kohn和Sham等从总交换能的角度出发%给})
{得}出$\alpha=2/3$\upcite{PR140-A1133_1965}，Slater从交换势的角度出发，%给})
{得}出$\alpha=1$\upcite{PR81-385_1951}，对原子、分子体系的计算表明，选取$\alpha$约为0.7\upcite{AQC6-1_1972,Slater-4_1974}结果更好。

对于用自旋密度泛函方法处理的自旋极化体系%的情况})
，%此时的})
交换能{泛函的}形式可以写成\upcite{PRA20-397_1979}
\begin{equation} \label{eq:dft-7}
E_X[\rho_{\alpha},\rho_{\beta}]=\frac12E_X^{0}[2\rho_{\alpha}]+\frac12E_X^{0}[2\rho_{\beta}]
\end{equation}
其中$E_X^{0}[\rho]=E_X\bigl[\dfrac 12\rho_{\alpha},\dfrac 12\rho_{\beta}\bigr]$

相关能包含相同自旋电子间的相互作用和不同自旋电子间相互作用，不可能把相关能泛函也分解成为两种自旋电子的相关能的贡献和。一般采用Stoll\upcite{TCA49-143_1978}的定义，将局域密度近似的相关能泛函写成如下形式：
\begin{equation}
  E_C=E_C^{\alpha\beta}+E_C^{\alpha\alpha}+E_C^{\beta\beta}
  \label{eq:dft-17}
\end{equation}

把局域密度近似引入自旋密度泛函理论就是通常所说的局域{自旋}密度近似(local spin-density approximation, LSDA)，LSDA中的交换能{泛函}%形})
{表达}式可以%结合})
{由}\eqref{eq:dft-7}式和\eqref{eq:dft-6}式%给})
{得}出{。}相关能泛函的形式可以写成：
$$E_C^{\mathrm{LSDA}}[\rho_{\alpha},\rho_{\beta}]=\int{\varepsilon_C(\rho,\varsigma)\rho\,\textrm{d}^3\vec{r}}$$
其中$\varsigma=(\rho_{\alpha}-\rho_{\beta})
/(\rho_{\alpha}+\rho_{\beta})
$，$\rho_{\alpha}$和$\rho_{\beta}$分别是自旋为$\alpha$和自旋为$\beta$的电子密度。均匀电子气相关能的解析形式很难求得，但是可以用数值方法计算得到%其})
数值结果，Vosko等\upcite{CJP58-1200_1980}用解析式拟合%出})
精确的数值计算结果，得到局域密度近似下的相关能泛函VWN表达式{。}%但是})
他们给出的表达式%太})
{比较}复杂，Perdew和Wang\upcite{PRB45-13244_1992}%又})
把它改%写})
成%为})
{比较}简单的形式，得到%的})
{一个电子的}相关能%密度})
{的}表达式为：
\begin{equation}
  \label{eq:dft-8}
  \varepsilon_C^{\mathrm{LDA}}(r_s,\varsigma)=\varepsilon_C^{\mathrm{LDA}}(r_s,0)+a_C(r_s)\dfrac{f(\varsigma)}{f^{\prime\prime}(\varsigma)}(1-\varsigma^4)+[\varepsilon_C(r_s,1)-\varepsilon_C(r_s,0)]f(\varsigma)\varsigma^4
\end{equation}
其中$r_s=\left[\dfrac 3{4\pi}\bigl(\rho_{\alpha}+\rho_{\beta}\bigr)\right]^{-1/3}$，$f(\varsigma)=[(1+\varsigma)^{4/3}+(1-\varsigma)^{4/3}-2]/(2^{4/3}-2)$，\\$\varepsilon_C(r_s,0)$，$\varepsilon_C(r_s,1)$和$-a_C(r_s)$由经验公式
$$G(r_s,A,\alpha_1,\beta_1,\beta_2,\beta_3,\beta_4,p)=-2A(1+\alpha_1r_s)\ln\left[1+\dfrac1{2A(\beta_1r_s^{1/2}+\beta_2r_s+\beta_3r_s^{3/2}+\beta_4r_s^{p+1})}\right]$$
计算{，}其中{$G$可以是$\varepsilon_C(r_s,0)$，$\varepsilon_C(r_s,1)$或$-a_C(r_s)$}，$A$，$\alpha_1$，$\beta_1$，$\beta_2$，$\beta_3$，$\beta_4$，$p$%都})
是{特定的}参数。

由于局域密度近似下的交换-相关能泛函形式是从均匀电子气模型中得到的，一般认为它对于电子密度变化比较小的体系结果比较好，%而})
{可能}不适用于电子密度剧烈变化的原子、分子和一些固体体系。%事实上})
{不过}，局域密度近似用于原子、分子特别是某些扩展体系仍能得到很合理的结果，其原因是%因为})
局域密度近似下的交换-相关能密度和交换-相关孔在每一点上的数值都来自一个体系的精确的值，因此满足交换-相关电荷总量为一个电子电荷的负值的约束条件，交换-相关能泛函的误差%在不同位置})
{能部分}相互抵消。

\subsubsection{梯度校正和GGA泛函}
%由于})
LDA是建立在均匀电子气模型基础上{的}，而原子、分子和很多扩展体系的电子密度远非均匀，%通常})
由LDA计算得到的分子的原子化能{通常}过大。%要})
{为}进一步提高计算精度，%就})
需要对电子密度的非均匀性加以校正。一般是通过在交换-相关能泛函中引入电子密度%的})
梯度来加以校正{，}即构造GGA泛函。可以将GGA交换能泛函写成一般形式：
\begin{equation}
	E_X^{\mathrm{GGA}}=E_X^{\mathrm{\mathrm{LDA}}}-\sum_{\sigma}\int F[x_{\sigma}]\rho_{\sigma}^{\frac43}(\vec r)\textrm{d}^3r{=E_X^{\mathrm{LDA}}+E_X^{\mathrm{NL}}}
  \label{eq:dft-18}
\end{equation}
其中$x_{\sigma}=|\nabla\rho_{\sigma}|\rho_{\sigma}^{-4/3}$称为约化梯度，是一个无量纲的量。%依据})
{按}所用$F${泛函}的形式不同，%目前})
{现有}GGA的交换能泛函可以分为两大类，一类以Becke\upcite{PRA38-3098_1988}在1988年提出的表达式为基础：
\begin{equation}
  \label{eq:dft-9}
  E_X^{\mathrm{NL}}=-b\sum_{\sigma}\int\rho_{\sigma}^{4/3}\dfrac{x_{\sigma}^2}{1+6bx_{\sigma}\sinh^{-1}x_{\sigma}}\textrm{d}^3r
\end{equation}
其中$b$\,=\,0.0042{。}%$X_{\sigma}=|\nabla\rho_{\sigma}|\rho_{\sigma}^{-4/3}$})
这是目前常用的交换能非局域校正公式。属于这一类的交换能泛函有PW91\upcite{PRB46-6671_1992,PRB48-4978_1993,PRB54-16533_1996,PRB57-14999_1998}、FT97\upcite{MP91-847_1997}、CAM(A)和CAM(B)\upcite{JCP99-8765_1993}等等。另一类泛函的$F${泛函}采用有理函数{形式，}包括幂函数和有理分式的泛函。属于这一类的交换能泛函有B86\upcite{JCP84-4524_1986}、P86$_x$\upcite{PRB33-8800_1986}、LG\upcite{PRA47-4681_1993}、PBE\upcite{PRL77-1396_1996,IBID78-1396_1997}、VSXC\upcite{JCP109-400_1998}等。B86采用的$F${泛函}的形式是
\begin{equation}
	F^{\mathrm{B86}}=\left[1+1.296\left(\dfrac{x_{\sigma}}{(24\pi^2)^{1/3}}\right)^2+14\left(\dfrac{x_{\sigma}}{(24\pi^2)^{1/3}}\right)^4+0.2\left(\dfrac{x_{\sigma}}{(24\pi^2)^{1/3}}\right)^6\right]^{1/15}
  \label{eq:dft-19}
\end{equation}

最常用的相关能非局域校正{$E_C^{\mathrm{NL}}$}是Perdew和Wang\upcite{PRB33-8822_1986}提出的%形})
{表达}式以及把$E_C^{\mathrm{NL}}$和$E_C^{\mathrm{LDA}}$合{并}在一起计算的LYP%形})
{表达}式\upcite{PRB37-785_1988}。Perdew和Wang的表达式为：
\begin{equation} 
  \label{eq:dft-10}
  E_C^{\mathrm{NL}}=\int{d^{-1}\textrm{e}^{-\varphi}C[\rho]|\nabla\rho|^2\rho^{-4/3}\textrm{d}^3r}
\end{equation}
其中$\varphi=1.745\times 0.11\times C[\infty]|\nabla\rho|/(C[\rho]\rho^{7/6})$，%\\%\linebreak

$d=2^{1/3}\left[\biggl(\dfrac {1+\varsigma}2\biggr)^{5/3}+\biggl(\dfrac{1-\varsigma}2\biggr)^{5/3}\right]^{1/2}$，

$C[\rho]=a+(b+ar_s+\beta r_s^2)(1+\gamma r_s+\delta r_s^2+10^4\beta r_s^3)^{-1}$，\\%\linebreak
$a$，$b$，$\alpha$，$\beta$，$\gamma$，$\delta$为通过拟合{实验数据}得到的参数。

LYP相关能泛函是%利用})
{在}Colle-Salvetti公式\upcite{TCA37-329_1975}{的基础上}导出的，其具体形式为：
\begin{equation}
  \begin{aligned}
	  E_C^{\mathrm{LYP}}=&-a\int\dfrac{\gamma(\vec r)}{1+d\rho^{-1/3}}\biggl\{\rho+2b\rho^{-5/3}\biggl[2^{2/3}C_F\rho_{\alpha}^{8/3}+2^{2/3}C_F\rho_{\beta}^{8/3}-\rho t_w\biggr.\biggr.\\
  &+\left.\left.\dfrac19\left(\rho_{\alpha}t_w^{\alpha}+\rho_{\beta}t_w^{\beta}\right)+\dfrac1{18}\left(\rho_{\alpha}\nabla^2\rho_{\alpha}+\rho_{\beta}\nabla^2\rho_{\beta}\right)\right]\textrm{e}^{-c\rho^{-1/3}}\right\}\textrm{d}\vec r
  \end{aligned}
  \label{eq:dft-20}
\end{equation}
其中$\gamma(\vec r)=2\left[1-\dfrac{\rho_{\alpha}^2(\vec r)+\rho_{\beta}^2(\vec r)}{\rho^2(\vec r)}\right]$，$t_w(\vec r)=\dfrac18\cdot\dfrac{|\nabla\rho(\vec r)|^2}{\rho(\vec r)}-\dfrac18\nabla^2\rho(\vec r)$，$C_F=\dfrac3{10}(3\pi^2)^{2/3}$，\\$a$，$b$，$c$，$d$%是})
{为通过拟合实验数据确定的}常数。

\subsubsection{\textit{meta}-GGA泛函}
由于密度泛函理论的交换-相关能中包含了实际体系与无相互作用参考体系动能{之}间的差别，%因此如果})
在交换和相关能泛函中包含动能密度$\tau$，%则})
计算结果可能会更合理一些。表达式中包含动能密度$\tau$为变量的泛函称为\textit{meta}-GGA泛函。动能密度的一般定义为
\begin{equation}
	\tau=\sum_i^{\mathrm{occ}}|\nabla\Psi_i|^2
  \label{eq:dft-21}
\end{equation}
其中$\Psi_i$是含自旋的轨道。

van~Voorhis和Scuseria于1998年提出了基于密度矩阵展开理论\upcite{PRC5-1472_1972,PRC11-1031_1975}的VSXC泛函\upcite{JCP109-400_1998}，得到的交换能密度泛函为：
\begin{equation}
  \begin{aligned}
    E_X[\rho_{\alpha},\rho_{\beta}]&=\sum_{\sigma}\int\rho_{\sigma}^{4/3}{f(x_{\sigma},z_{\sigma})
    }\textrm{d}\vec r\\
    f(x_{\sigma},z_{\sigma})&=\left(\dfrac a{\Gamma_{\sigma}(x,z)}+\dfrac{bx_{\sigma}^2+cz_{\sigma}}{\Gamma_{\sigma}^2(x,z)}+\dfrac{dx_{\sigma}^4+ex_{\sigma}^2z_{\sigma}+f_0z_{\sigma}^2}{\Gamma_{\sigma}^3(x,z)}\right) 
  \end{aligned}
  \label{ed:def-22}
\end{equation}
其中$x_{\sigma}=|\nabla\rho_{\sigma}|\rho_{\sigma}^{-4/3}$，$z_{\sigma}=\tau\rho_{\sigma}^{-5/3}-C_F$，$C_F=\dfrac35(3\pi^2)^{2/3}$，$\Gamma_{\sigma}(x,z)=1+\alpha(x_{\sigma}^2+z_{\sigma})$，\linebreak $\alpha$为常数；$a$，$b$，$c$，$d$，$e$，$f_0$均为参数。

{在}相关能泛函的计算中，将相同自旋电子的相关作用和不同自旋{电子}的相关作用分开%计算})
{处理}\upcite{TCA49-143_1978}，%有})
\begin{equation}
  \begin{aligned}
    E_C=&E_C^{\alpha\beta}+E_C^{\alpha\alpha}+E_C^{\beta\beta}\\
    E_C^{\sigma\sigma^{\prime}}=&\int f^{\sigma\sigma^{\prime}}(x,z)e_{c\sigma\sigma^{\prime}}^{\mathrm{LDA}}\textrm{d}\vec r\\
    E_C^{\sigma\sigma}=&\int f^{\sigma\sigma}(x_{\sigma},z_{\sigma})
    D_{\sigma}e_{c\sigma\sigma}^{\mathrm{LDA}}\textrm{d}\vec r
  \end{aligned}
  \label{eq:dft-23}
\end{equation}
其中$x^2=x_{\alpha}^2+x_{\beta}^2$，$z=z_{\alpha}+z_{\beta}$，{$e_C^{\mathrm{LDA}}$为相关能密度，}$D_{\sigma}$是相同自旋自相互作用的校正系数，表达式为：
$$D_{\sigma}=1-\dfrac{x_{\sigma}^2}{4(z_{\sigma}+C_F)}$$

1989年Becke和Roussel从类氢原子的密度矩阵表达式出发，提出%了})
BR89交换能泛函\upcite{PRA39-3761_1989}。
交换势为：
$$U_{X\sigma}=-4\pi\int_0^{\infty}\rho_X^{\mathrm{H}}(a,b,s)s\textrm{d}s=-\left(1-\textrm{e}^{-x}-\frac12x\textrm{e}^{-x}\right)/b$$
其中$\rho_X^{\mathrm{H}}(a,b,s)$是类氢原子的交换孔函数，%$x$})
{$x=ab$}，由{逐点求解}超越方程
$$\dfrac{x\textrm{e}^{-2/3}}{x-2}=\dfrac23\pi^{2/3}\dfrac{\rho_{\sigma}^{5/3}}{Q_{\sigma}}$$
%逐点求解})
得到。方程中$Q_{\sigma}$\,=\,$\dfrac16\left[\nabla^2\rho_{\sigma}-2\gamma\bigl(\tau_{\sigma}-\dfrac{(\nabla\rho_{\sigma})^2}{4\rho_{\sigma}}\bigr)\right]$，其中$\gamma$是参数，$\tau_{\sigma}$为动能密度{，}\linebreak $b$%根据})
\,=\,$\left[\dfrac{x^3\textrm{e}^{-x}}{8\pi\rho_{\sigma}}\right]^{1/3}$%计算得到})
{，$a=x/b$。}%最后得到})
总的交换能为：$E_{X\sigma}=\dfrac12\displaystyle\int\rho_{\sigma}U_{X\sigma}\textrm{d}\vec r$

对于分子体系，交换孔函数比较离域，BR89模型往往高估交换能密度，Becke用LDA动能与精确动能密度的比值为参数，提出了B00修正模型\upcite{JCP112-4020_2000}。修正后的交换能泛函形式为：
\begin{equation}
	E_{X\sigma}^{\mathrm{BR89}}=\int e_{X\sigma}^{\mathrm{BR89}}[1+f_{X\sigma}(t_{\sigma})
  ]\textrm{d}\vec r
  \label{eq:dft-24}
\end{equation}
其中$f_{X\sigma}$是校正函数，$e_{X\sigma}^{\mathrm{BR89}}$是BR89模型中的交换能密度。$f_{X\sigma}(t_{\sigma})=a_t(\omega_{\sigma}-2\omega_{\sigma}^3+\omega_{\sigma}^5)$，$\omega_{\sigma}=\dfrac{t_{\sigma}-1}{t_{\sigma}+1}$，$t_{\sigma}=\dfrac{\tau_{\sigma}^{\mathrm{LSDA}}}{\tau_{\sigma}^{\mathrm{exact}}}$，式中$a_t$为经验参数。

%若})
加上{电子}相关{能}作用(采用Bc88相关能泛函\upcite{JCP88-1053_1988})
，得到的交换-相关能{泛函的}表达式为：
\begin{equation}
	E_{XC}=E_X^{\mathrm{BR89}}+c_{\mathrm{opp}}E_{C\mathrm{opp}}^{\mathrm{Bc88}}+c_{\mathrm{par}}E_{C\mathrm{par}}^{\mathrm{Bc88}}+a_{mix}\sum_{\sigma}\int e_x^{\mathrm{BR89}}f_{x\sigma}(t_{\tau})\textrm{d}\vec r
  \label{eq:dft-25}
\end{equation}
其中$E_{C\mathrm{opp}}^{\mathrm{Bc88}}$和$E_{C\mathrm{par}}^{\mathrm{Bc88}}$分别是自旋平行、反平行的相关能。拟合得到的参数$a_{mix}$的优选值为0.135。

\subsubsection{{杂化(hybrid)}泛函}
%还有})
{目前}一种很常用的交换-相关能{泛函}是%采用})
把Hartree-Fock%的相关})
{交换}能%形式})
与近似%的})
交换-相关能%相})
混合%的方法})
\upcite{JCP104-1040_1996,JCP98-5648_1993}，其表达式%可以写})
为：
\begin{equation}
	E_{XC}^{\mathrm{hybrid}}=\alpha E_X+(1-\alpha)\bar{E}_X+\bar{E}_C
 \label{eq:dft-11}
\end{equation}
其中$E_X$是Hartree-Fock%的相关})
{交换}能%形式})
，$\bar{E}_X$和$\bar{E}_C$分别为近似的交换能和相关能泛函，通常$\alpha$取值为0.28左右。%例如})
“半对半(half and half)”泛函\upcite{JCP98-1372_1993}%，其})
{的}表达式%可以写成})
{为}：
\begin{equation}
	E_{XC}^{\mathrm{HH}}=\dfrac12E_{X}^{\mathrm{HF}}+\dfrac12E_{XC}^{LSDA}
  \label{eq:dft-26}
\end{equation}

杂化型的密度泛函方法可以统一地归为一类，称为绝热关联方法\textrm{(Adiabatic Connection Method, ACM)}。所谓绝热关联，就是定义以$\lambda$为参量的Hamiltonian
\begin{equation}
  \hat H_{\lambda}=-\dfrac12\sum_i\nabla_i^2+\sum_iV_i(\rho,\lambda)+\dfrac{\lambda}2\sum_{i\neq j}\dfrac1{r_{ij}}
  \label{eq:dft-27}
\end{equation}
其中$V_i(\rho,\lambda)$是与$\lambda$有关的外势，选择$V_i(\rho,\lambda)$使不同$\lambda$时的密度$\rho$保持恒定。%显然})
当$\lambda$=0时对应于无相互作用参考体系；而当$\lambda$=1时对应%着})
真实体系。当$\lambda$连续地从0变化到1时，无相互作用参考体系绝热地过渡到真实体系。利用一级微扰理论\upcite{JCP88-1053_1988}，可以得到交换-相关能%地})
{的}表达式为：
$$E_{XC}=\dfrac12\int\rho(\vec r)d\vec r\int\dfrac{\bar h_{xc}(\vec r,\vec r\,^{\prime})}{|\vec r-\vec r\,^{\prime}|}\textrm{d}\vec r\,^{\prime}$$
其中$$\bar h_{xc}(\vec r,\vec r\,^{\prime})=\int_0^1h_{xc}^{\lambda}(\vec r,\vec r\,^{\prime})\textrm{d}\lambda$$
$h_{xc}^{\lambda}(\vec r,\vec r\,^{\prime})$是与$\hat H_{\lambda}$相应的交换-相关孔函数。令
$$U_{XC}^{\lambda}=\dfrac12\int\rho(\vec r)d\vec r\int\dfrac{h_{xc}^{\lambda}(\vec r,\vec r\,^{\prime})}{|\vec r-\vec r\,^{\prime}|}\textrm{d}\vec r\,^{\prime}$$
于是有$$E_{XC}=\int_0^1U_{XC}^{\lambda}\textrm{d}\lambda$$
如果%选用最})
简单%的})
{地}假设，%认为})
$U_{XC}^{\lambda}$与$\lambda$满足线性关系，就%可以})
得到
$$E_{XC}=\dfrac12U_{XC}^0+\dfrac12U_{XC}^1=\dfrac12E_X^{HF}+\dfrac12U_{XC}^1$$
这就是“半对半(half and half)”泛函以及%所有})
{其它}杂化型泛函的理论基础。

属于这一类的泛函还有B3P\upcite{JCP98-5648_1993}、B3LYP\upcite{JPC98-11623_1994}、B1B95\upcite{JCP104-1040_1995}、B97\upcite{JCP107-8554_1997}、B98\upcite{JCP109-2092_1998}、PBE0\upcite{JCP110-6158_1999}~等。其中最流行的B3LYP泛函的具体形式为：
$$E_{XC}^{\mathrm{B3LYP}}=(1-a)E_X^{\mathrm{LSDA}}+aE_X^{\mathrm{HF}}+bE_X^{\mathrm{B88}}+cE_C^{\mathrm{LYP}}+(1-c)E_C^{\mathrm{LSDA}}$$
其中$a$，$b$，$c$%均})
为参数。

%一般来说})
这几种泛函的近似公式的计算结果{一般}差别并不很大，对各种近似泛函的计算结果可以参考文献\inlinecite{JCP98-5612_1993,CPL316-160_2000}。

\subsubsection{绝热关联与交换-相关孔}
为了进一步讨论近似交换-相关泛函与精确交换-相关泛函的差别，按照绝热关联\textrm{(adiabatic connection)}过程，引入电子间相互作用:
\begin{equation}
	V_{\mathrm{C},\lambda}=\sum_{ij}\dfrac{\lambda}{|\vec r_i-\vec r^{\prime}_j|}=\lambda V_{\mathrm{C}}
	\label{eq:electron-electron}
\end{equation}
于是，对于外势下的均匀电子气的\textrm{Hamiltonian}表示
\begin{equation}
	H_{\lambda}=H_{0,\lambda}+\sum_iV_{\mathrm{ext}}(\vec r_i)=(T+V_{\mathrm{C},\lambda})+\sum_iV_{\mathrm{ext}}(\vec r_i)
	\label{eq:H-interaction}
\end{equation}
要求当体系密度$\rho(\vec r)$不变时，\textrm{Hamiltonian}的外势部分$\sum\limits_iV_{\mathrm{ext}}(\vec r_i)$对体系能量的贡献不变。不难看出$\lambda=0$时，表示无相互作用电子体系，$\lambda=1$表示完全相互作用电子体系。对于给定的电荷密度，最小化\textrm{Hamiltonian}的体系能量
\begin{equation}
	E_{\lambda}[\rho]=\min_{\psi|\rho}\langle\Psi|H_{0,\lambda}|\Psi\rangle=\int V_{\mathrm{ext}}(\vec r)\rho(\vec r)\mathrm{d}^3r=F_{\lambda}[\rho]+\int V_{\mathrm{ext}}(\vec r)\rho(\vec r)\mathrm{d}^3 r
	\label{eq:min-E}
\end{equation}
式中$F_{\lambda}[\rho]$定义为
\begin{equation}
	F_{\lambda}[\rho]=\min_{\psi|\rho}\langle\Psi|H_{0,\lambda}|\Psi\rangle
	\label{eq:General-Functional-lambda}
\end{equation}
体系能量最小化过程中，要求对于给定的电荷密度$\rho$，存在对应的体系波函数$\Psi$。根据\textrm{Hellmann-Feynman}定理\upcite{Hellmann_Book,PR56-340_1939}，当\textrm{Hamiltonian}仅依赖于参数$\lambda$时，可有对含$\lambda$的\textrm{Schr\"odinger}方程
\begin{equation}
	H_{\lambda}|\psi_{\lambda}\rangle=E_{\lambda}|\psi_{\lambda}\rangle
	\label{eq:lambda-Schrodinger}
\end{equation}
对$\lambda$求微分可有:
\begin{equation}
	H_{\lambda}^{\prime}|\psi_{\lambda}\rangle+H_{\lambda}|\psi_{\lambda}^{\prime}\rangle=E_{\lambda}^{\prime}|\psi_{\lambda}\rangle+E_{\lambda}|\psi_{\lambda}^{\prime}\rangle
	\label{eq:Hellmann-Feynman}
\end{equation}
$H_{\lambda}^{\prime}$、$\psi_{\lambda}^{\prime}$和$E_{\lambda}^{\prime}$表示对$\lambda$的导数。方程两边左乘$\langle\Psi_{\lambda}|$，并利用式\eqref{eq:lambda-Schrodinger}，可有
\begin{equation}
	\dfrac{\mathrm{d}E_{\lambda}}{\mathrm{d}\lambda}=\dfrac{\langle\psi_{\lambda}|\mathrm{d}H_{\lambda}/\mathrm{d}\lambda|\psi_{\lambda}\rangle}{\langle\psi_{\lambda}|\psi_{\lambda}\rangle}
	\label{eq:HF-QMD}
\end{equation}
当$|\psi_{\lambda}\rangle$对应$F_{\lambda}$的变分基态($H_{0,\lambda}$的最低能级的本征态)，则参照\textrm{Hellmann-Feynman}定理，可以将$F_{\lambda}$的导数写成
\begin{equation}
	\dfrac{\mathrm{d}F_{\lambda}}{\mathrm{d}\lambda}=\langle\psi_{\lambda}|V_\mathrm{C}|\psi_{\lambda}\rangle
	\label{eq:HF-F}
\end{equation}
因此对于完全相互作用体系的通用泛函
\begin{equation}
	F_{\lambda=1}[\rho]=T_{\mathrm{K-S}}[\rho]+\int_0^1\langle\psi_{\lambda}|V_{\mathrm{C}}|\psi_{\lambda}\rangle\mathrm{d}\lambda
	\label{eq:General-F}
\end{equation}
因此交换-相关能可以表示为
\begin{equation}
	\begin{aligned}
		E_{XC}=&F_{\lambda=1}[\rho]-T_{\mathrm{K-S}}[\rho]-\sum\dfrac12\int\rho(\vec r)\dfrac1{|\vec r-\vec r^{\prime}|}\rho(\vec r^{\prime})\mathrm{d}^3\vec r\mathrm{d}^3\vec r^{\prime}\\
		=&\int_0^1\langle\psi_{\lambda}|V_{\mathrm{C}}|\psi_{\lambda}\rangle\mathrm{d}\lambda-E_{\mathrm{Hartree}}
	\end{aligned}
	\label{eq:EX-delta}
\end{equation}
上述推导不难看出，相互作用体系的动能部分贡献的交换-相关，可以通过绝热关联\textrm{(Adiabatic Connection)}近似 ，计算\textrm{Coulomb}项对参数$\lambda$积分贡献得到。

根据密度矩阵理论，对于电子体系，电子1和2分别位于位置$x$和$x^{\prime}$的概率密度可以表示为
\begin{equation}
	P(x,x^{\prime})=\int|\Psi(x,x^{\prime},x_3,\cdots,x_N)|^2\mathrm{d}x_3\cdots\mathrm{d}x_N
	\label{eq:density-matrix}
\end{equation}
单电子密度则可以表示为
\begin{equation}
	\rho(x)=N\int|\Psi(x,x^{\prime},x_3,\cdots,x_N))|^2\mathrm{d}x^{\prime}\mathrm{d}x_3\cdots\mathrm{d}x_N
	\label{eq:single-density}
\end{equation}
并有
\begin{equation}
	\int\rho(x)\mathrm{d}x=N
	\label{eq:N-totlal}
\end{equation}
不难看出，存在等式
\begin{equation}
	N\int P(x,x^{\prime})\mathrm{d}x^{\prime}=\rho(x)
	\label{eq:relation-matrix-N-tot}
\end{equation}
\textrm{Coulomb}相互作用能用密度矩阵表示
\begin{equation}
	E_{\mathrm{C}}=\dfrac{N(N-1)}2\int P(x,x^{\prime})\dfrac1{|\vec r-\vec r^{\prime}|}\mathrm{d}x\mathrm{d}x^{\prime}
	\label{eq:Coulomb-matrix}
\end{equation}
引入交换-相关孔$n_{XC}(x,x^{\prime})$，满足
\begin{equation}
	N(N-1)P(x,x^{\prime})=\rho(x)\rho(x^{\prime})+\rho(x)n_{XC}(x,x^{\prime})
	\label{eq:ex-ch-hole}
\end{equation}
这样定义的交换-相关孔，表明当考虑电子间相互作用后，第一个电子位于位置$x$时，根据体系的平均密度分布给出的第二个电子的概率分布。因此可以有
\begin{equation}
	E_{\mathrm{C}}=E_{\mathrm{Hartree}}+\dfrac12\int \rho(x)n_{XC}(x,x^{\prime})\dfrac1{|\vec r-\vec r^{\prime}|}\mathrm{d}x\mathrm{d}x^{\prime}
	\label{eq:Coulomb-Hartree-XC}
\end{equation}
式\eqref{eq:Coulomb-Hartree-XC}的第二项，即交换-相关能来自对交换-相关孔的\textrm{Coulomb}积分形式。

通过对交换-相关孔性质的讨论，可以获得精确交换-相关能需要满足的某些条件。由式\eqref{eq:N-totlal}和波函数的归一化条件，可有
\begin{equation}
	\int P(x,x^{\prime})\mathrm{d}x\mathrm{d}x^{\prime}=1
	\label{eq:normalisation}
\end{equation}
对交换-相关孔定义式\eqref{eq:ex-ch-hole}两边对变量$x^{\prime}$(即空间坐标$\vec r^{\prime}$和自旋$s$)的积分，可有
\begin{equation}
	(N-1)\rho(x)=N\rho(x)+\rho(x)\int n_{XC}(x,x^{\prime})\mathrm{d}x^{\prime}
	\label{eq:ex-ch-hole-integrate}
\end{equation}
不难得出，交换-相关孔应满足条件
\begin{equation}
	\int n_{XC}(x,x^{\prime})\mathrm{d}x^{\prime}=-1
	\label{eq:ex-ch-hole-condition}
\end{equation}
因此式\eqref{eq:Coulomb-Hartree-XC}的第二项可以视为交换-相关效应对经典\textrm{Coulomb}相互作用(\textrm{Hartree}能)的修正。换言之，交换-相关孔可以视为具有单位正电荷的一种电荷密度分布。如果将交换-相关孔分为交换孔和相关孔两部分之和
\begin{equation}
	n_{XC}(x,x^{\prime})=n_X(x,x^{\prime})+n_C(x,x^{\prime})
	\label{eq:ex-ch-devode}
\end{equation}
根据\textrm{Hartree-Fock}方法，波函数采用\textrm{Slater}行列式形式，即可完全地考虑交换作用，由此确定精确的交换孔满足
\begin{equation}
	\int n_{X}(x,x^{\prime})\mathrm{d}x^{\prime}=-1
	\label{eq:ex-hole-condition}
\end{equation}
不难看出，精确相关孔满足
\begin{equation}
	\int n_{C}(x,x^{\prime})\mathrm{d}x^{\prime}=0
	\label{eq:ch-hole-condition}
\end{equation}

由于原子中的电子密度的不均匀性比分子中的不均匀性更强，交换-相关能的梯度校正部分对原子{计算}影响比较大，%能够})
比较明显%的})
{地}降低原子的能量，而对分子{计算}影响较小{。}与LDA的结果相比，交换-相关能的梯度校正使得分子的原子化能{计算值}降低，键长{计算值}增加，总体%上得到的})
{看计算}结果%也})
有明显改善。%也正是这些})
精度较高的交换-相关能泛函%形式})
的提出，使%得})
密度泛函理论%能够广泛应用于定量研究})
{在}原子、分子、团簇和固体等%各种})
体系的%性质。})
{定量研究中得到广泛的应用并且得到满意的结果。}

有必要指出的是，目前使用的{近似能量}泛函%形式})
还存在问题，比较集中的有：
\begin{enumerate}
	\item 自相互作用抵消不干净；
	\item 对于简并基态和近简并基态还存在%严重问题})
{难于合理解决的困难}；
	\item 对于弱相互作用体系，如van der Waals相互作用不能忽略时，{计算结果}还有很大的误差。%由于})
\end{enumerate}
电子与{其}自身没有相互作用，因此交换-相关能中包含的电子自相互作用部分应该和电子间Coulomb能中包含的自相互作用部分完全抵消，而目前使用的{近似能量密度}泛函%形式})
都不%具有})
{严格满足}这个%特点})
{条件}，因此{用}密度泛函方法计算电子数很少的体系一般都会有比较大的误差。{在体系处于}简并基态和近简并基态的%问题一直是密度泛函理论比较困难的部分})
{情况下}，当%体系的})
基态电子密度%取自})
{用}不同的简并轨道{计算}时%密度泛函方法得到的})
体系能量应该%保持一致})
{不变}，%现在的})
{而现有近似能量密度}泛函%形式也})
不具有这个%特点})
{性质}，%由于多数原子的基态是简并的，})
用%密度泛函方法})
{来}计算%的})
原子基态能量就%会})
随%着选})
取%的})
不同简并原子轨道{计算电子密度}而不同\upcite{PRB43-6865_1991,CPL265-481_1997}，%这})
给计算分子体系的原子化能造成很大%的})
困难。密度泛函方法计算弱相互作用体系出现%问题})
的{困难}主要%原因})
是弱相互作用体系%通常键})
{中的相互作用}能和现%在所用的近似泛函形式})
{有近似能量密度泛函}本身的{计算}误差%可以相比})
{在同一量级}。%但是})
寻找精度%有})
明显提高的近似能量密度泛函%形式})
是难度很大的工作，目前常用的{近似能量密度}泛函%形式})
主要是二十世纪八十年代后期和九十年代初提出的，此后虽然%也有一些})
{提出过很多}近似{能量密度}泛函%形式被提出来})
，%但是})
与以前提出的%泛函形式})
相比%都})
{精度并}没有很明显的%改进})
{提高}。
\section{固体能带理论}
%\subsection{\rm{Blochl}定理}
%\subsection{能带与$\vec k$空间}
%\section{从赝势到\rm{PAW}方法}
%\subsection{赝势方法}
%\subsection{\rm{APW}方法与\rm{PAW}方法}
%\subsection{\rm{VASP}中的{\rm{PAW}}方法}
%\section{一些专题}
%\subsection{$k$空间积分与总能量}
%\subsection{第一原理分子动力学与\rm{Car-Parrinello}方案}
%\subsection{电子相关的解决方案:~\rm{DFT}+$U$和\rm{GW}近似}
%\subsection{自旋-轨道耦合}
%\subsection{电子的激发态和线性响应}
固体能带理论是凝聚态物理中最成功的理论之一，是固体电子理论的支柱。上世纪二十年代末和三十年代初，量子力学的运动规律确立后，在应用量子力学研究金属电导理论的过程中发展起来了固体能带理论，最初的成就在于定性的阐述了晶体中电子运动的普遍性特点。固体的很多性质，如振动光谱、磁有序、电导率、光学介电函数等，原则上都可以由固体能带理论阐述和解释。上世纪五十年代以后，特别是六十年代，由于研究固体的实验工作的重大发展，提供了大量的实验数据；随着大型高速计算机的应用，使得能带理论的研究从定性的普遍规律性发展到对具体的材料复杂的能带结构计算。

固体能带理论是一个近似理论，它的主要任务是确定固体的电子能级，也就是能带。在固体中存在着大量的电子，它们的运动是相互关联的，每个电子的运动都与其他电子的运动相关，对于这种多电子体系严格求解是不可能的。能带理论是单电子近似理论，就是把每个电子的运动看成是独立的在一个由其他电子和原子核构成的等效势场中的运动。在大多数情况下，人们最关心的是价电子，在原子构成分子或者固体的过程中，价电子的的运动状态往往发生很大的变化。而内层电子的的变化则比较小，因此可以将原子核与内层电子近似看成是一个离子实，这样价电子的等效势场，包括离子实的势场，其他价电子的平均势场以及考虑电子波函数的反对称性而引入的交换作用。单电子近似最早用于研究多电子原子，又称Hatree-Fock自洽场方法。

能带理论的出发点是固体中的电子不再束缚于个别的原子，而是在整个固体范围内运动，称为“共有化电子”。在讨论共有化电子的运动状态的时候，假定原子实处于平衡位置，而把原子实偏离平衡位置的影响看成微扰。对于理想晶体，原子排列成晶格，晶格具有平移周期性，因此势能函数也应具有平移周期性。晶体中的单电子在一个具有晶格周期性的等效势场中运动，其波动方程是：
\begin{equation}\label{eq:solid-1}
  \biggl[-\dfrac{\hbar^2}{2m}\nabla^2+V(\vec r)\biggr]\psi=E\psi
\end{equation}
其中势能$V(\vec r)$满足：
\begin{equation}\label{eq:solid-2}
  V(\vec r)=V(\vec r+\vec R_n)
\end{equation}
$\vec R_n$为任意格矢。

\section{平移周期性和倒格矢}
理想晶体中的原子(离子实)的排列具有平移对称性，%即晶体的周期性规则的重复结构，可以看作是一个或者一组原子(或离子实)以某一特定的方式在空间中作周期性平移的结果。
这是理想晶体和普通分子的最大区别。晶体中被平移的重复单元，称为基元(basis)，不同的晶体有着不同的基元；晶体中重复单元的平移方式，一般以抽象的空间点阵表示，称为晶格(crystal lattice)。基元以相同的方式排列在晶格的结点上。理想晶体是基元和晶格的统称。


晶格通常用Bravais格子的形式表示。Bravais格子是矢量
\begin{equation}
  \vec R_n=n_1\vec{\mathbf a}_1+n_2\vec{\mathbf a}_2+n_3\vec{\mathbf a}_3
  \label{eq:def-Bravais}
\end{equation}
全部端点的集合，其中$n_1,n_2,n_2$为整数，$\vec{\mathbf a}_1,\vec{\mathbf a}_2,\vec{\mathbf a}_3$为三个不共面的矢量，称Bravais格子的基矢(primitive vector)，$\vec R_n$称为Bravais格子的格矢(Bravais lattice vector, BLV)，其端点称为格点(lattice site)。所有格点的环境相同，在几何上是完全等价的。Bravais格子是一个无穷延伸的完全理想点阵，反映了晶体中原子(或离子实)的周期性排列。%，是对具有平移对称性的实际晶体合理的理想抽象。


晶体中体积最小的周期性重复单元称为原胞(primitive cell)。%将原胞平移所有Bravais格子可能的格矢$\vec R_n$，将精确的占据理想晶体的所有空间，既不会有遗漏，也不会有重复。
习惯上将原胞取为以基矢为棱边的平行六面体，体积为：
\begin{equation}
  \Omega_0=\vec{\mathbf a}_1\cdot(\vec{\mathbf a}_2\times\vec{\mathbf a}_3)
  \label{eq:vol_pri_lattice}
\end{equation}
基矢的取法并不唯一。%，因此原胞有多种取法，但是任意选取的原胞并不一定具有Bravais格子的全部点群对称性。
为保证原胞具有Bravais格子的全部点群对称性，通常选用Wigner-Seitz(WS)原胞。以晶格中某一格点为中心，作其与最近邻格点连线的垂直平分面，所有这些平面围成的以该格点为中心的最小体积称为属于该点的WS原胞。实际上WS原胞是Bravais的对称化原胞。

%体积V内$\vec R_n$的格点，对于$\vec r$处的格点密度的贡献为$\delta(\vec r-\vec R_n)$，于是
%\begin{equation}
%  \rho(\vec r)=\sum_{\vec R_n\in V}\delta(\vec r-\vec R_n)
%\end{equation}
%其中$\vec R_n$为Bravais格子的格矢。

%对于忽略表面效应的理想晶体，Bravais格子满足平移对称性要求
%\begin{equation}
%  \rho(\vec r+\vec R_n)=\rho(\vec r)
%  \label{eq:solid-3}
%\end{equation}
%对所有属于BLV的$\vec R_n$成立。即$\rho(\vec r)$为周期函数，平移Bravais格子的任意格矢保持不变。除了$\rho(\vec r)$之外，晶体的其他一些性质，如质量密度，电子云密度，离子实产生的势场等也都是周期函数。一般可以写成：
Bravais格子中的周期函数一般可以写成：
\begin{equation}
  F(\vec r+\vec R_n)=F(\vec r)
  \label{eq:solid-4}
\end{equation}
%对于所属的BLV的$\vec R_n$成立。将$F(\vec r)$展开成为Fourier级数
%\begin{equation}
%  F(\vec r)=\sum_g A(\vec g)e^{i\vec g\cdot\vec r}
%  \label{eq:solid-5}
%\end{equation}
%其中系数
%\begin{equation}
%  A(\vec g)=\dfrac1{\Omega}\int_{\Omega}F(\vec r)e^{-i\vec g\cdot\vec r}d\vec r=\dfrac{1}{\Omega}\int_{\Omega}F(\vec r+\vec R_n)e^{-i\vec g\cdot\vec r}d\vec r
%  \label{eq:solid-6}
%\end{equation}
%令$\vec r'=\vec r+\vec R_n$，则式\eqref{eq:solid-5}化为：
%\begin{equation}
%  A(\vec g)=\dfrac{1}{\Omega}\int_{\Omega}F(\vec r')e^{-i\vec g\cdot\vec r'}d\vec r'\cdot e^{i\vec g\cdot\vec R_n}=A(\vec g)e^{i\vec g\cdot\vec R_n}
%  \label{eq:solid-7}
%\end{equation}
%即
%\begin{equation}
%  A(\vec g)[1-e^{i\vec g\cdot\vec R_n}]=0
%  \label{eq:solid-8}
%\end{equation}
%对所有的$\vec g$，如果$A(\vec g)=0$，则$F(\vec r)\equiv0$，因此要求$e^{i\vec g\cdot\vec R_n}=1$。

对Bravais格子中所有的格矢$\vec R_n$，满足
\begin{equation}
  e^{i\vec G_h\cdot\vec R_n}=1
  \label{eq:def_recip_1}
\end{equation}
即
\begin{equation}
  \vec G_h\cdot\vec R_n=2\pi m, \quad m\mbox{为整数}
  \label{eq:def_recip_2}
\end{equation}
的全部$\vec G_h$端点，构成该Bravais格子(即正格子(direct lattice))的倒格子(reciprocal lattice)。%将$\vec R_n=n_1\vec{\mathbf a}_1+n_2\vec{\mathbf a}_2+n_3\vec{\mathbf a}_3$代入上式，有
%\begin{equation}
%  n_1\vec G_h\cdot\vec{\mathbf a}_1+n_2\vec G_h\cdot\vec{\mathbf a}_2+n_3\vec G_h\cdot\vec{\mathbf a}_3=2\pi m
%  \label{eq:solid-9}
%\end{equation}
%对任意整数$n_1$、$n_2$和$n_3$都成立，因此要求
%\begin{equation}
%  \vec G_h\cdot\vec{\mathbf a}_1=2\pi h_1,\quad\vec G_h\cdot\vec{\mathbf a}_2=2\pi h_2,\quad\vec G_h\cdot\vec r_3=2\pi h_3,\quad h_1,h_2,h_3\mbox{为整数}
%  \label{eq:solid-10}
%\end{equation}
因此倒格矢可以写成
\begin{equation}
  \vec G_h=h_1\vec{\mathbf b}_1+h_2\vec{\mathbf b}_2+h_3\vec{\mathbf b}_3,\quad h_1,h_2,h_3\mbox{为整数。}
   \label{eq:solid-11}
\end{equation}
且满足
\begin{equation}
   \vec{\mathbf b}_i\cdot\vec{\mathbf a}_j=2\pi\delta_{ij}
  \label{eq:solid-12}
\end{equation}
%因为$\vec{\mathbf a}_1,\vec{\mathbf a}_2,\vec{\mathbf a}_3$不在同一平面内，条件所确定的
倒格矢$\vec G_h$垂直于Miller指数为($h_1,h_2,h_3$)的晶面系， 晶面间距为$d=\dfrac{2\pi}{|\vec G_h|}$。倒空间(reciprocal space)的基矢$\vec{\mathbf b}_1,\vec{\mathbf b}_2,\vec{\mathbf b}_3$可以表示为：
\begin{equation} \label{eq:solid-13}
\left\{ \begin{aligned}
  \vec{\mathbf b}_1&=2\pi \dfrac{\vec{\mathbf a}_2\times\vec{\mathbf a}_3}{\vec{\mathbf a}_1\cdot(\vec{\mathbf a}_2\times\vec{\mathbf a}_3)} \\
  \vec{\mathbf b}_2&=2\pi \dfrac{\vec{\mathbf a}_3\times\vec{\mathbf a}_1}{\vec{\mathbf a}_1\cdot(\vec{\mathbf a}_2\times\vec{\mathbf a}_3)} \\
  \vec{\mathbf b}_3&=2\pi \dfrac{\vec{\mathbf a}_1\times\vec{\mathbf a}_2}{\vec{\mathbf a}_1\cdot(\vec{\mathbf a}_2\times\vec{\mathbf a}_3)} \\
\end{aligned} \right.
\end{equation}
倒格矢是倒空间中的Bravais格子。%可以证明，
倒格子的原胞体积$\Omega^{\ast}$与正格子的原胞体积成反比。即
\begin{equation}
  \Omega^{\ast}=\vec{\mathbf b}_1\cdot(\vec{\mathbf b}_2\times\vec{\mathbf b}_3)=\dfrac{(2\pi)^3}{\Omega_0}
  \label{eq:solid-14}
\end{equation}
类似地，倒格子空间中的WS原胞称为第一Brillouin区(first Brillouin zone)。同一晶格的正格子和倒格子有相同的点群对称性，由于WS原胞具有Bravais的全部点群对称性，所以第一Brillouin区也具有晶格点群的全部对称性。


%与Bravais格子有相同平移对称性的物理量的Fourier展开中，只存在与格矢为倒格矢的分量，其他分量的系数为零。即
%满足\eqref{eq:solid-4}的函数$F(\vec r)$的Fourier展开式为
%\begin{equation}
%  \begin{split}
%  F(\vec r)&=\sum_{\vec G_h}A(\vec G_h)e^{i\vec G_h\cdot\vec r},\\
%  \mbox{其中}\hspace{1\ccwd} A(\vec G_h)&=\dfrac1{\Omega_0}\int_{\Omega_0}F(\vec r)e^{-i\vec G_h\cdot\vec r}d\vec r
%  \label{eq:solid-15}
%  \end{split}
%\end{equation}


%根据正格子与倒格子的关系，可以证明；

\section{Bloch定理}
\subsection{Bloch定理及其证明}
{\bf bloch定理:}\hspace{1\ccwd}{\it 对于具有周期性的势场，单电子的Schr\"odinger方程为}
\begin{equation}
  \psi(\vec r+\vec R_n)=e^{i\vec k\cdot\vec R_n}\psi(\vec r)
  \label{eq:bloch}
\end{equation}
{\it其中$\vec k$为一矢量}。式\eqref{eq:bloch}表明当平移晶格矢量$\vec R_n$时，波函数只增加了相位因子$e^{i\vec k\cdot\vec R_n}$。式\eqref{eq:bloch}就是Bloch定理\cite{ZP52-555_1928}。根据Bloch定理，波函数可以写成
\begin{equation}
  \psi(\vec r)=e^{i\vec k\cdot\vec r}u(\vec r)
  \label{eq:blochpw}
\end{equation}
其中$u(\vec r)$具有与晶格同样的周期性，即$u(\vec r+\vec R_n)=u(\vec r)$。
%\begin{equation}
%  u(\vec r+\vec R_n)=u(\vec r)
%  \label{eq:solid-16}
%\end{equation}
式\eqref{eq:blochpw}称为Bloch波函数，它是平面波与周期函数的乘积。Bloch定理的简单证明如下：


周期性势场反映了晶体的平移对称性，即晶格平移任意格矢$\vec R_n$后势场是不会改变的。引入平移算符$\hat T_{\vec R_n}$，其中$\vec R_n$是Bravais格子的任一格矢，定义为$\hat T_{\vec R_n}$作用于任意函数$f(\vec r)$上，使得矢量$\vec r$平移$\vec R_n$，即
\begin{equation}
  \hat T_{\vec R_n}f(\vec r)=f(\vec r+\vec R_n)
  \label{eq:solid-17}
\end{equation}


在晶体中单电子运动的Hamiltonian\,$\hat H=-\dfrac{\hbar^2}{2m}\nabla^2+V(\vec r)$
%\begin{equation}
%  \hat H=-\dfrac{\hbar^2}{2m}\nabla^2+V(\vec r)
%  \label{eq:solid-18}
%\end{equation}
具有晶格周期性，而
\begin{equation}
  \begin{split}
  \hat T_{\vec R_n}\hat Hf(\vec r)&=\biggl[-\dfrac{\hbar^2}{2m}\nabla_{\vec r +\vec R_n}^2+V(\vec r+\vec R_n)\biggr]f(\vec r+\vec R_n)\\
  &=\biggl[-\dfrac{\hbar^2}{2m}\nabla_{\vec r}^2+V(\vec r)\biggr]f(\vec r+\vec R_n)\\
  &=\hat H\hat T_{\vec R_n}f(\vec r)
  \end{split}
  \label{eq:solid-19}
\end{equation}
%因为其中$\nabla_{\vec r+\vec R_n}$只是相应的对$\partial/\partial x$，$\partial/\partial y$，$\partial/\partial z$中的变量$x$，$y$，$z$改变一个常数，不会影响微分算符(即微分算符与平移操作无关)。可见晶体Hamiltonian式\eqref{eq:solid-18}具有平移周期性$\hat H(\vec r)=\hat H(\vec r+\vec R_n)$，而且
因此$\hat T_{\vec R_n}$与$\hat H$是对易的，即$\hat T_{\vec R_n}\hat H=\hat H\hat T_{\vec R_n}$。
%\begin{equation}
%  \hat T_{\vec R_n}\hat H=\hat H\hat T_{\vec R_n}
%  \label{eq:solid-20}
%\end{equation}
%根据量子力学一般原理，两个对易算符具有相同的本征函数。因此，
可以选择$\hat H$的本征态，使它同时为各平移算符$\hat T_{\vec R_n}$的本征态。即有
\begin{equation} \label{eq:solid-21}
\left\{ \begin{aligned}
  \hat H\psi(\vec r)&=E\psi(\vec r) \\
  \hat T_{\vec R_n}\psi(\vec r)&=\lambda_{\vec R_n}\psi(\vec r) \\
\end{aligned} \right.
\end{equation}
$\lambda_{\vec R_n}$是对应的平移算符的本征值。根据平移算符的定义\eqref{eq:solid-17}，
\begin{equation}
  \psi(\vec r+\vec R_n)=\lambda_{\vec R_n}\psi(\vec r)
  \label{eq:solid-22}
\end{equation}
%为了确定本征值$\lambda_{\vec R_n}$的值，
由波函数的归一化条件得，%$\int_{\Omega_0}|\psi(\vec r)|^2d\vec r=\int_{\Omega_0}|\psi(\vec r+\vec R_n)|^2d\vec r=1$
%\begin{equation}
%  \int_{\Omega_0}|\psi(\vec r)|^2d\vec r=\int_{\Omega_0}|\psi(\vec r+\vec R_n)|^2d\vec r=1
%  \label{eq:solid-23}
%\end{equation}
$|\lambda_{\vec R_n}|^2=1$。
%\begin{equation}
%  |\lambda_{\vec r}|^2=1
%  \label{eq:solid-24}
%\end{equation}
于是$\lambda_{\vec R_n}$可以写成$\lambda_{\vec R_n}=e^{i\beta\vec R_n}$
%\begin{equation}
%  \lambda_{\vec R_n}=e^{i\beta\vec R_n}
%  \label{eq:solid-25}
%\end{equation}
的形式，即$\psi(\vec r+\vec R_n)$和$\psi(\vec r)$仅差了一个相位因子。


此外，考虑到任意两次相继的平移$\vec R_n$和$\vec r_m$，相当于一次平移$\vec R_n+\vec r_m$之和，即满足：
\begin{equation}
  \begin{aligned}
   \hat T_{\vec R_n}\hat T_{\vec r_m}\psi&=\hat T_{\vec R_n}\lambda_{\vec r_m}\psi=\lambda_{\vec r_m}\lambda_{\vec R_n}\psi \\
   \hat T_{\vec R_n}\hat T_{\vec r_m}\psi&=\hat T_{\vec R_n+\vec r_m}\psi=\lambda_{\vec R_n+\vec r_m}\psi
  \end{aligned}
  \label{eq:solid-26}
\end{equation}
因此平移算符的本征值$\lambda_{\vec R_n}$必须满足关系$\lambda_{\vec R_n+\vec r_m}=\lambda_{\vec r_m}\lambda_{\vec R_n}$。
%\begin{equation}
%  \lambda_{\vec R_n+\vec r_m}=\lambda_{\vec r_m}\lambda_{\vec R_n}
%  \label{eq:solid-27}
%\end{equation}
%将式\eqref{eq:solid-25}代入，两边取对数，有
%\begin{equation}
%  \beta_{\vec R_n+\vec r_m}=\beta_{\vec R_n}+\beta_{\vec r_m}
%  \label{eq:solid-28}
%\end{equation}
%式\eqref{eq:solid-28}仅当$\beta$与$\vec R_n$之间呈线性关系才能满足，即$\beta_{\vec R_n}=\vec k\cdot\vec R_n$，
有$\lambda_{\vec R_n}=e^{i\vec k\cdot\vec R_n}$。
%\begin{equation}
%  \lambda_{\vec R_n}=e^{i\vec k\cdot\vec R_n}
%  \label{eq:solid-29}
%\end{equation}
$\vec k$称为简约波矢。


只要$\hat H$具有平移对称性，对于任意Bravais格子的格矢$\vec R_n$，其本征函数满足
\begin{equation}
  \hat T_{\vec R_n}\psi(\vec r)=\psi(\vec r+\vec R_n)=\lambda_{\vec R_n}\psi(\vec r)＝e^{i\vec k\cdot\vec R_n}\psi(\vec r)
  \label{eq:solid-30}
\end{equation}
这就是Bloch定理。


\subsection{简约波矢$\vec k$的取值和物理意义}
为了确定波矢$\vec k$的值，必须引入边界条件。选择周期性边条件(也称Born-Von Karman条件)。
\begin{equation}
  \left\{
  \begin{aligned}
     \psi(\vec r+N_1\vec{\mathbf a}_1)&=\psi(\vec r)\\
     \psi(\vec r+N_2\vec{\mathbf a}_2)&=\psi(\vec r)\\
     \psi(\vec r+N_3\vec{\mathbf a}_3)&=\psi(\vec r)\\
   \end{aligned}\right.
  \label{eq:solid-31}
\end{equation}
其中$\vec{\mathbf a}_i(i=1,2,3)$是Bravais格子的三个基矢，$N_1$、$N_2$、$N_3$分别是沿着基矢$\vec{\mathbf a}_1$、$\vec{\mathbf a}_2$、$\vec{\mathbf a}_3$方向的原胞数，$N=N_1N_2N_3$是晶体中的原胞总数。因此$\vec k$(相应的本征值$\lambda_{\vec R_n}$)的取值将受到限制。根据Bloch定理，%将式\eqref{eq:bloch}代入\eqref{eq:solid-31}得
\begin{equation}
  \psi(\vec r+N_i\vec{\mathbf a}_i)=e^{iN_i\vec k\cdot\vec{\mathbf a}_i}\psi(\vec r),\quad i=1,2,3
  \label{eq:solid-32}
\end{equation}
这就要求$e^{iN_i\vec k\cdot\vec{\mathbf a}_i}=1$，
%\begin{equation}
%  e^{iN_i\vec k\cdot\vec{\mathbf a}_i}=1,\quad i=1,2,3
%  \label{eq:solid-33}
%\end{equation}
%或者等价地
%\begin{equation}
%  N_i\vec k\cdot\vec{\mathbf a}_i=2\pi l_i,\quad l_i\mbox{为整数},\quad i=1,2,3
%  \label{eq:solid-34}
%\end{equation}
将波矢$\vec k$用相应的倒格子的基矢$\vec{\mathbf b}_i(i=1,2,3)$表示，%即
%\begin{equation}
%  \vec k=k_1\vec{\mathbf b}_1+k_2\vec{\mathbf b}_2+k_3\vec{\mathbf b}_3
%  \label{eq:solid-35}
%\end{equation}
%代入式\eqref{eq:solid-34}，并
利用正交关系式\eqref{eq:solid-12}$\vec{\mathbf a}_i\cdot\vec{\mathbf b}_j=2\pi\delta_{ij}$，%有
%\begin{equation}
%  \vec k=\dfrac{l_1}{N_1}\vec{\mathbf b}_1+\dfrac{l_2}{N_2}\vec{\mathbf b}_2+\dfrac{l_3}{N_3}\vec{\mathbf b}_3
%  \label{eq:solid-36}
%\end{equation}
许可的简约波矢$\vec k$可以看成是倒格子空间中以$\dfrac{\vec{\mathbf b}_i}{N_i}(i=1,2,3)$为基矢的Bravais格子的格矢。


每个许可的$\vec k$值由上述Bravais格子的格点表示，在$\vec k$空间内所占的体积
\begin{equation}
  \Delta\vec k=\dfrac{\vec{\mathbf b}_1}{N_1}\cdot\biggl(\dfrac{\vec{\mathbf b}_2}{N_2}\times\dfrac{\vec{\mathbf b}_3}{N_3}\biggr)=\dfrac1N\vec{\mathbf b}_1\cdot(\vec{\mathbf b}_2\times\vec{\mathbf b}_3)
  \label{eq:solid-37}
\end{equation}
由于$\vec{\mathbf b}_1\cdot(\vec{\mathbf b}_2\times\vec{\mathbf b}_3)$是倒格子原胞的体积，因此倒格子空间原胞许可的$\vec k$的数目等于实空间中晶体的总的原胞数目。


倒格子原胞体积为$\dfrac{(2\pi)^3}{\Omega_0}$，%式\eqref{eq:solid-14}，
$\Omega_0$是正格子的原胞体积，$N\Omega_0=V$，因此$\vec k$空间中许可态的态密度
\begin{equation}
  \dfrac1{\Delta\vec k}=\dfrac V{8\pi^3}
  \label{eq:solid-38}
\end{equation}
对于用平面波描述的自由电子，$\hbar\vec k$是电子的动量\cite{Yanshousheng}。但是对于Bloch电子，简约波矢$\vec k$并不比例于电子的动量。动量算符$p=-i\hbar\nabla$作用于Bloch波函数%(式\eqref{eq:blochpw})上，
\begin{equation}
  \begin{split}
  -i\hbar\nabla\psi_{\vec k}&=-i\hbar\nabla(e^{i\vec k\cdot\vec r}u_{\vec k}(\vec r))\\
  &=\hbar\vec k\psi_{\vec k}-i\hbar e^{i\vec k\cdot\vec r}\nabla u_{\vec k}(\vec r)
  \end{split}
  \label{eq:solid-39}
\end{equation}
并不能写成一个简单的常数乘以$\psi_{\vec k}$，因此$\psi_{\vec k}$也不是动量算符的本征函数。


简约波矢$\vec k$是对应于平移操作本征值的量子数，它的物理意义是表示原胞之间电子波函数位相的变化。不同$\vec k$值表明原胞间的位相是不同的。但是需要指出的是，如果$\vec k$改变一个倒格子矢量
\begin{displaymath}
  \vec G_n=n_1\vec{\mathbf b}_1+n_2\vec{\mathbf b}_2+n_3\vec{\mathbf b}_3,\quad(n_1,n_2,n_3\mbox{为整数})
\end{displaymath}
完全不会影响本征值$\lambda_{\vec R_n}$。因此，为了使$\vec k$能与平移算符的本征值$\lambda_{\vec R_n}$一一对应，必须要把$\vec k$的取值限制在一定的范围之内，使得它既能概括所有不同的本征值$\lambda_{\vec R_n}$的取值，同时又没有任意两个$\vec k$相差一个倒格矢$\vec G_n$。最方便的办法是将$\vec k$限制在$\vec k$空间的WS原胞中。

\section{近自由电子近似和紧束缚近似}
近自由电子近似和紧束缚近似是能带理论中的两个最基本的理论模型。近自由电子近似是在自由电子气模型的基础上引入弱的周期性势场，使得电子的运动发生变化，将弱的周期性势场视为微扰，从而引入能带和能隙的概念。近自由电子近似对于有较多{\it s}\,电子和{\it p}\,电子的金属，是很好的近似。紧束缚近似从原子有序排列形成晶体过程出发，考察近邻原子的波函数叠加而引起束缚于特定原子的单电子态的运动变化，可以得到电子的原子能级与晶体能带的相互联系。其物理图像和结果更适用于过渡金属的{\it d}\,电子和{\it f}\,电子以及固体中的其他内层电子。

\subsection{周期场的近自由电子近似}
近自由电子近似假定周期性变化的势场起伏比较小，将势场的平均值$\bar V$代替$V(\vec r)$作为零级近似，将周期性起伏$[V(\vec r)-\bar
V]$作为微扰处理。以一维周期体系为例，根据微扰理论，得到一级修正的波函数为：
\begin{equation}
  \begin{aligned}
    \psi_k&=\psi_k^0+\psi_k^{(1)} \\
    &=\dfrac1{\sqrt L}e^{ikx}+\sum_n\dfrac{V_n}{\dfrac{\hbar^2}{2m}\left[k^2-\left(k+\dfrac na2\pi\right)^2\right]}\cdot\dfrac1{\sqrt V}e^{i(k+2\pi\frac na)x} \\
    &=\dfrac1{\sqrt L}e^{ikx}\biggl\{1+\sum_n\dfrac{V_n}{\dfrac{\hbar^2}{2m}\left[k^2-\left(k+\dfrac na2\pi\right)^2\right]}e^{i2\pi\frac nax}\biggr\}
  \end{aligned}
  \label{eq:solid-40}
\end{equation}
式中$L=Na$为一维晶格长度，$N$为原胞数目，$a$为晶格常数(原子间距)。在$x$改变为$a$的整数倍的时候，求和项不变，说明括号内为一周期函数，符合Bolch函数的形式，可以写成一个自由粒子的波函数乘上具有晶格周期性的函数。


二级能量微扰写成
\begin{equation}
  E_k^{(2)}=\sum_n\dfrac{|V_n|^2}{\dfrac{\hbar^2}{2m}\left[k^2-\left(k+\dfrac na2\pi\right)^2\right]}
  \label{eq:solid-41}
\end{equation}
注意到当
\begin{equation}
  k^2=\left(k+\dfrac na2\pi\right)^2
  \label{eq:solid-42}
\end{equation}
即当$k$是$\dfrac{\pi}a$的整数倍($k=-\dfrac{n\pi}a,n\mbox{为任意整数}$)时微扰理论失败，此时必须采用简并态微扰理论，
\begin{equation}
  E_{\pm}=\dfrac{\hbar^2k^2}{2m}\pm|V_n|
  \label{eq:solid-43}
\end{equation}
这样，周期性的起伏势场使得自由电子在波矢$k=\dfrac12G_n=\dfrac{\pi}an$处断开，能量突变为$2|V_n|$。这种断开使得准连续的电子能谱出现能隙(energy gap)。在能隙范围内没有许可的电子态，电子能级分裂成为一系列的能带。


一维体系的情形很容易推广到二维和三维情形。如果从$\vec k$空间中的原点出发，沿特定方向考察$E(\vec k)$的变化，在近自由电子近似情形下，$E(\vec k)$将像自由电子一样，比例于$k^2$呈抛物线形式变化。在跨越Brillouin区边界或者Bragg平面是一倒格矢$\vec G_h$的垂直平分面，能量跃迁的大小约为$2|V_{\vec G_h}|$的量级。$V_{\vec G_h}$是近自由电子近似Fourier展开中与$\vec G_h$联系的项的系数。

\subsection{紧束缚近似}
紧束缚近似的出发点是，电子在一个原子附近时，将主要受到该原子场的作用，将其他原子场的作用看作微扰作用，据此可以得到电子的原子能级与晶体中能带之间的相互联系。


如果完全不考虑原子间相互影响，在某格点$\vec r_m=m_1\vec{\mathbf a}_1+m_2\vec{\mathbf a}_2+m_3\vec{\mathbf a}_3$
%\begin{displaymath}
%  \vec r_m=m_1\vec{\mathbf a}_1+m_2\vec{\mathbf a}_2+m_3\vec{\mathbf a}_3
%\end{displaymath}
附近的电子将以原子束缚态$\varphi_i(\vec r-\vec R_m)$的形式环绕$\vec R_m$点运动，(这里假定简单晶格，每个原胞含有一个原子)，$\varphi_i$表示孤立原子波动方程的本征态
\begin{equation}
  \left[-\dfrac{\hbar^2}{2m}\nabla^2+V(\vec r-\vec r_m)\right]\varphi_i(\vec r-\vec r_m)=\varepsilon_i\varphi_i(\vec r-\vec r_m)
  \label{eq:solid-44}
\end{equation}
$V(\vec r-\vec r_m)$是$\vec R_m$格点的原子势场，$\varepsilon_i$为某个原子能级。晶体中的电子波动方程为
\begin{equation}
   \left[-\dfrac{\hbar^2}{2m}\nabla^2+U(\vec r)\right]\psi_i(\vec r)=E\psi_i(\vec r)
  \label{eq:solid-45}
\end{equation}
$U(\vec r)$为周期性势场，它是各格点原子势场之和。在紧束缚近似中将式\eqref{eq:solid-44}看成零级近似，将$U(\vec r)-V(\vec r-\vec R_m)$看成微扰。环绕不同的格点，将有$N$个类似的波函数，它们具有相同的能量$\varepsilon_i$，即$N$重简并态。这实际上是把原子间相互影响看作微扰的简并微扰方法，微扰后的状态是$N$个简并态的线性组合，即用原子轨道$\varphi_i(\vec r-\vec R_m)$的线性组合来构造晶体中电子共有化运动的轨道$\psi(\vec r)$，因而也称为原子轨道线性组组合法(Linear Combination of Atomic Orbitals, LCAO)。因此有
\begin{equation}
  \psi(\vec r)=\sum_ma_m\varphi_i(\vec r-\vec R_m)
  \label{eq:solid-46}
\end{equation}
%将式\eqref{eq:solid-46}代入晶体中电子的波动方程\eqref{eq:solid-45}，利用式\eqref{eq:solid-44}得到：
%\begin{equation}
%  \sum_ma_m[\varepsilon_i+U(\vec r)-V(\vec r-\vec r_m)]\varphi_i(\vec r-\vec r_m)=E\sum_ma_m\varphi_i(\vec r-\vec r_m)
%  \label{eq:solid-47}
%\end{equation}
当原子间距比原子轨道半径大时，不同格点的$\varphi_i$的重叠很小，将近似认为
\begin{displaymath}
  \int\varphi_i^{\ast}(\vec r-\vec R_m)\varphi_i(\vec r-\vec R_n)d\vec r=\delta_{nm}
\end{displaymath}

根据紧束缚近似，可以看出原子能级和能带的对应关系可以分为下列几种情形：

1. 对于最简单的情况，一个原子能级$\varepsilon_i$对应一个能带，原子的各个不同能级，在固体中将产生一系列相应的能带。愈低的能带愈窄，愈高的能带愈宽。这是由于能量最低的能带对应于最内层的电子，它们的电子轨道“半径”很小，在不同原子间很少相互重叠，因此能带很窄。能量较高的外层电子轨道，在不同的原子间将有较大的重叠，从而形成较宽的带。只有在简单情况下，原子能级与能带间有简单的对应关系。

2. 对复杂的情形，原子能级与能带间并不存在简单的一一对应关系。在形成晶体过程中，不同原子态之间可能相互混合。通常可以用能带宽度反映混合的程度。对内层电子，能带宽度较小，能级与能带间有简单的一一对应关系；外层电子能带较宽，能带与能带间对应关系变得复杂，可以认为主要由几个能级相近的原子态相互组合形成能带，略去其它较多原子态的影响。

3. 对于复式晶格，每个原胞内有若干个原子，可以认为原胞中各原子之间先形成分子轨道，能带与分子轨道之间有相互对应的关系。

在紧束缚近似中，能带中的电子波函数可以写成原子波函数的Bloch和：
\begin{displaymath}
  \psi_{\vec k}^i=\dfrac1{\sqrt N}\sum_ne^{i\vec k\cdot\vec R_n}\varphi_i(\vec r-\vec R_n)
\end{displaymath}
对于任何能带，Bloch函数都可以写成类似的形式：
\begin{equation}
  \psi_{n\vec k}=\dfrac1{\sqrt N}\sum_ne^{i\vec k\cdot\vec R_n}W_n(\vec r-\vec R_n)
  \label{eq:solid-48}
\end{equation}
其中$V_n(\vec r-\vec R_n)$称为Wannier函数。
\begin{displaymath}
  W_n(\vec r-\vec R_n)=\dfrac1{\sqrt N}\sum_{\vec k}e^{-i\vec k\cdot\vec R_n}\psi_{n\vec k}
\end{displaymath}
也就是说，一个能带的Wannier函数是由同一能带的Bloch函数所定义的。Wannier函数之间完全正交，%容易证明
%\begin{displaymath}
%  \int W_n^{\ast}(\vec r-\vec R_n)W_n(\vec r-\vec R_{m'})d\vec r=\delta_{mm'}
%\end{displaymath}
因此，Bloch函数的集合与Wannier函数的集合是两组完备的正交函数集，它们之间由幺正矩阵相联系。


在紧束缚近似中，如果近似忽略原子波函数的交叠，近似的有
\begin{displaymath}
  \int \varphi_i^{\ast}(\vec r-\vec R_n)\varphi_i(\vec r-\vec R_m)d\vec r=\delta_{nm}
\end{displaymath}
在这种情况下，Wannier函数就是各个格点上孤立原子的波函数。如果某些能带与紧束缚近似模型相差很远，这时Wannier函数是很少保留孤立原子波函数的信息，但它仍然比较定域化。在讨论那些电子空间局域性起重要作用问题时，Wannier函数将是比较好的工具。
